% Compile this single file twice with pdflatex (or use latexmk -pdf).
% All mathematical content and the bibliography are included below.
\documentclass[11pt,a4paper]{article}
\usepackage[T1]{fontenc}
\usepackage[utf8]{inputenc}
\usepackage{lmodern,microtype}
\usepackage[left=25mm,right=25mm,top=24mm,bottom=24mm,headheight=15pt,headsep=6mm]{geometry}
\usepackage{amsmath,amssymb,amsthm,mathtools}
\usepackage{booktabs,array,longtable,tabularx}
\usepackage[dvipsnames]{xcolor}
\usepackage{fancyhdr,enumitem}
\usepackage{hyperref,bookmark}
\definecolor{ink}{HTML}{163E46}
\definecolor{muted}{HTML}{46656B}
\hypersetup{colorlinks=true,linkcolor=ink,citecolor=ink,urlcolor=ink,
pdftitle={One-zero proofs of irrationality and transcendence},pdfsubject={A catalogue of 100 arithmetic conclusions},pdfauthor={}}
\urlstyle{same}
\newtheorem{theorem}{Theorem}[section]
\newtheorem{lemma}[theorem]{Lemma}
\newtheorem{proposition}[theorem]{Proposition}
\newtheorem{corollary}[theorem]{Corollary}
\theoremstyle{definition}
\newtheorem{definition}[theorem]{Definition}
\theoremstyle{remark}
\newtheorem{remark}[theorem]{Remark}
\setlist{nosep,topsep=4pt}
\setlength{\emergencystretch}{1.7em}
\setlength{\parindent}{1.15em}
\setcounter{tocdepth}{1}
\setcounter{secnumdepth}{2}
\pagestyle{fancy}
\fancyhf{}
\fancyhead[L]{\small\sffamily\color{muted}One-zero proofs of irrationality and transcendence}
\fancyhead[R]{\small\sffamily\color{muted}\thepage}
\renewcommand{\headrulewidth}{0.25pt}
\raggedbottom
\begin{document}
\thispagestyle{plain}
\noindent{\footnotesize\sffamily\color{muted} RESEARCH MANUSCRIPT\hfill 15 November 2026}\par
\vspace{1.3em}
\noindent{\fontsize{26}{30}\selectfont\bfseries\color{ink} One-zero proofs of irrationality\\[1pt] and transcendence\par}
\vspace{.55em}
\noindent{\Large A catalogue of one hundred arithmetic conclusions}\par
\vspace{1.0em}
\noindent\rule{\textwidth}{0.5pt}\par
\vspace{.6em}
\noindent\textbf{Abstract.}
A zero or a stationary point can determine an arithmetic constant when its defining equation eliminates every auxiliary function value. If the argument is transcendental, every nonconstant rational function of it with algebraic coefficients is transcendental. We develop this observation for core thresholds in random graphs and hypergraphs, including the positive solution of $e^t=1+t+t^2$ that controls the birth of a graph's $3$-core. For analytic quantiles, elliptic geometry, lattice walks and generating functions, specialization and independence theorems supply alternative arithmetic inputs. Positive linear forms and nonperiodic digits cover examples with irrationality conclusions. We give proofs for one hundred specified quantities, trace forty-three earlier research notes to the relevant arguments, and identify what remains to be proved in several proposed applications to classical constants.
\par\vspace{1em}
\noindent\textbf{Catalogue.} Each identifier names one quantity. The columns report the strongest of the two arithmetic conclusions established here, with stronger independence statements specified in the text. All decimal values are illustrative.
\begin{center}\small
\begin{tabularx}{\linewidth}{@{}lXrr@{}}\toprule
Entries & Objects & Transcendental & Irrational only\\\midrule
G01--G25 & Cores, allocation thresholds and probability laws & 25&0\\
A01--A25 & Quantiles, special functions and geometry & 23&2\\
M01--M25 & Modular forms, walks and growth rates & 23&2\\
N01--N25 & Series, digital constants and continued fractions & 7&18\\\midrule
& Total & \textbf{78}&\textbf{22}\\\bottomrule
\end{tabularx}
\end{center}
\vspace{.2em}
\noindent The count includes individual specializations of family theorems and quantities sharing a parameter. Each numbered application is proved where it is defined. Pairwise distinctness or algebraic independence across the whole catalogue is not assumed. The manuscript combines classical results, consequences of published theorems, and derivations from the earlier notes. It makes no priority claim for the one hundred entries.
\par\vspace{.85em}
\noindent\textbf{Contents}\par\smallskip
\noindent\begin{tabularx}{\textwidth}{@{}Xr@{}}
\hyperref[sec:method]{1. The one-zero method} & \pageref{sec:method}\\
\hyperref[sec:graph]{2. Core thresholds, allocations, and probability laws} & \pageref{sec:graph}\\
\hyperref[sec:analytic]{3. Quantiles, diffraction, and elliptic geometry} & \pageref{sec:analytic}\\
\hyperref[sec:mcat]{4. Modular values, lattice walks and growth rates} & \pageref{sec:mcat}\\
\hyperref[sec:arith]{5. Arithmetic series, digital constants, and effective avoidance} & \pageref{sec:arith}\\
\hyperref[sec:open]{6. Recovery identities and unresolved arithmetic questions} & \pageref{sec:open}\\
\hyperref[sec:archive]{Appendix A. Concordance with the forty-three earlier entries} & \pageref{sec:archive}\\
\hyperref[sec:directory]{Appendix B. Index of the one hundred quantities} & \pageref{sec:directory}\\
\hyperref[sec:ref]{References} & \pageref{sec:ref}
\end{tabularx}
\newpage
\section{The one-zero method}\label{sec:method}
Let $\overline{\mathbb Q}$ denote the field of all algebraic numbers in $\mathbb C$. Suppose a target $x$ is defined analytically, for instance as the value of an objective at its unique maximum. Introducing its maximizing argument $T$ helps if one can first prove an arithmetic property of $T$, then recover $x$ by an exact relation.
\begin{lemma}[Descent through a rational function]\label{lem:descent}
Let $T$ be transcendental and $R=P/Q\in\overline{\mathbb Q}(X)$ nonconstant. Then $R(T)$ is defined and transcendental. More generally, let $F\in\overline{\mathbb Q}[X,Y]$ satisfy $F(T,x)=0$. If $F(X,a)$ is not the zero polynomial for any algebraic $a$ under consideration, then $x$ is transcendental.
\end{lemma}
\begin{proof}
The nonzero polynomial $Q$ cannot vanish at $T$. If $a=P(T)/Q(T)$ were algebraic, $P-aQ$ would have algebraic coefficients and vanish at $T$. It is nonzero because $P/Q$ is nonconstant. The field of algebraic numbers is algebraically closed, giving the contradiction. For the last statement use $F(X,a)$ in place of $P-aQ$.
\end{proof}
The specialization condition matters. For example $(Y-1)(X-Y)=0$ holds at $(T,1)$ for every $T$. A reliable application either gives an explicit rational expression or checks its specialized polynomial. If one knows only $T\notin K$ for a field $K$, a nonconstant rational function can return to $K$; $(\sqrt2)^2=2$ illustrates the failure. An invertible fractional linear transformation over $K$ preserves exclusion from $K$, by applying its inverse.
\begin{proposition}[An exponential equation at the point]\label{prop:mechanism}
Let $A,B\in\overline{\mathbb Q}(X)$ be defined at $T$. If $A(T)\ne0$ and $e^{A(T)}=B(T)$, then $T$ is transcendental. In particular, every $x=R(T)$ with $R$ as in Lemma~\ref{lem:descent} is transcendental.
\end{proposition}
\begin{proof}
An algebraic $T$ would give a nonzero algebraic exponent $A(T)$ and an algebraic exponential $B(T)$. Hermite--Lindemann excludes this possibility \cite{aSiegel}. Apply Lemma~\ref{lem:descent} to $x$.
\end{proof}
\paragraph{From a model to an arithmetic assertion.}
Specify the observable and its units. Locate its defining point on a particular branch, by monotonicity, a sign change, a variational characterization, or an existing theorem. Evaluate the stationarity or zero equation before estimating anything numerically. For example, the otherwise transcendental factor $e^T$ in a Poisson tail becomes a polynomial in $T$ at the graph core threshold. One must still show that the rational expression left for the target is nonconstant. For complex special functions, a valid branch and the absence of poles enter at this step.
\paragraph{Other arithmetic inputs.}
If the point were algebraic, its function values might have to satisfy a relation excluded by a specialization theorem. Sections~\ref{sec:analytic} and~\ref{sec:mcat} state the hypotheses enabling this step for the functions used there. Alternatively, a set of values can generate a field of known transcendence degree. If a field of degree of transcendence at least $r$ is algebraic over $\overline{\mathbb Q}(x_1,\ldots,x_r)$, the $x_i$ are algebraically independent. This controls triples at a modular maximum and a reflected $q$-Stieltjes zero. Section~\ref{sec:arith} also treats sums whose defining series directly supplies a proof of irrationality.
\paragraph{What a recovery formula leaves to prove.}
In the proposed application $x=\gamma$, Lemma~\ref{lem:descent} would settle transcendence as soon as one has both an independently certified transcendental $T$ and a nonconstant $R\in\overline{\mathbb Q}(X)$ with $\gamma=R(T)$. A formula involving an additional uncontrolled value such as $Y_0'(T)$ does not supply this last assertion. Section~\ref{sec:open} gives the exact Bessel identities and locates the missing arithmetic step.


\section{Core thresholds, allocations, and probability laws}
\label{sec:graph}
Two random models can use the same Poisson tail and still have different
thresholds. We keep the numbers of choices, required degrees, and available
slots explicit. All 25 entries G01--G25 in this section are transcendental.
The tables give decimal approximations for identification. The exact
definitions in the theorems determine the numbers.

\subsection{The appearance of a core}
For integers $j\geq 0$ write $S_j(t)=\sum_{i=0}^j t^i/i!$ and put
$Q_h(t)=1-e^{-t}S_{h-1}(t)=\Pr\{\operatorname{Poisson}(t)\ge h\}$
for $h\ge1$. The $k$-core is the largest subgraph of minimum degree at
least $k$. For a $d$-uniform hypergraph we count incident hyperedges as degrees.
Molloy's threshold theorem \cite{gMolloy} gives the critical
mean degree as
\[
   b_{d,k}=\min_{t>0}\frac{t}{Q_{k-1}(t)^{d-1}},
   \qquad d,k\ge2,\quad(d,k)\ne(2,2).
\]
With $n$ vertices and $m$ uniformly chosen edges the edge density $m/n$ at
the threshold is $b_{d,k}/d$. For ordinary graphs, Pittel, Spencer and
Wormald \cite{gPSW} also identify $Q_k(T)$ as the limiting proportion of
vertices in the newborn $k$-core, where $T$ is the minimizing argument.

\par\hypertarget{cat-G01}{}\label{page:G01}\hypertarget{cat-G02}{}\label{page:G02}\hypertarget{cat-G03}{}\label{page:G03}\hypertarget{cat-G04}{}\label{page:G04}\hypertarget{cat-G05}{}\label{page:G05}\hypertarget{cat-G06}{}\label{page:G06}\hypertarget{cat-G07}{}\label{page:G07}\hypertarget{cat-G08}{}\label{page:G08}\hypertarget{cat-G09}{}\label{page:G09}\hypertarget{cat-G10}{}\label{page:G10}
\begin{theorem}\label{thm:gcore}
For every pair $(d,k)$ in this range, the minimizing argument $T$, the
critical mean degree $b_{d,k}$, and the value $Q_k(T)$ are transcendental.
\end{theorem}
\begin{proof}
Let $s=k-2$, $a=d-1$, and $B=a/s!$. The objective tends to infinity
both at zero and at infinity, since $a(s+1)>1$.
Using $Q_{s+1}'(t)=e^{-t}t^s/s!$, its stationary equation becomes
\[
 e^T=P(T),\qquad P(X)=S_s(X)+BX^{s+1}.
\]
It has exactly one positive solution. Indeed
$(e^t-S_s(t))/t^{s+1}$ is strictly increasing from $1/(s+1)!$ to
infinity on the positive axis, and $B>1/(s+1)!$.
The Hermite--Lindemann theorem makes $T$ transcendental. At this point
\[
 b_{d,k}=T\left(\frac{P(T)}{BT^{s+1}}\right)^a,\qquad
 Q_k(T)=\frac{(B-1/(s+1)!)T^{s+1}}{P(T)} .
\]
The first rational function grows like $X$ at infinity. The second
vanishes at zero and takes a positive value at $T$. Both are nonconstant,
so Lemma~\ref{lem:descent} applies.
\end{proof}

\paragraph{The original 3-core calculation.}
When $d=2$ and $k=3$, stationarity says $e^T=1+T+T^2$.
In this case the two observables reduce to
\[
   b_{2,3}=T+1+T^{-1},\qquad Q_3(T)=\frac{T^2}{2(1+T+T^2)} .
\]
Thus $3.35091887151167\ldots$ is the average-degree threshold and
$0.26758065498789\ldots$ is the asymptotic fraction of vertices present
at birth. These have different statistical meanings despite being
rational functions of the same argument.

\begin{center}\small
\begin{tabular}{@{}lll@{}}\toprule
ID & Observable & Approximation\\\midrule
G01 & Graph $3$-core threshold $b_{2,3}$ & 3.350918872\\
G02 & Graph $4$-core threshold $b_{2,4}$ & 5.149402747\\
G03 & Graph $5$-core threshold $b_{2,5}$ & 6.799275489\\
G04 & Graph $6$-core threshold $b_{2,6}$ & 8.365340770\\
G05 & Fraction in the newborn graph $3$-core & 0.267580655\\
G06 & Fraction in the newborn graph $4$-core & 0.438061713\\
G07 & Fraction in the newborn graph $5$-core & 0.538433562\\
G08 & Edge density $b_{3,2}/3$ for the $3$-uniform $2$-core & 0.818469161\\
G09 & Edge density $b_{4,2}/4$ for the $4$-uniform $2$-core & 0.772279840\\
G10 & Edge density $b_{3,3}/3$ for the $3$-uniform $3$-core & 1.552829940\\
\bottomrule\end{tabular}
\end{center}

\begin{corollary}[Fixed algebraic mean degree]\label{cor:gfixed}
Let $c>0$ be algebraic, $a=d-1$, $s=k-2$. For any real solution
$0<z<1$ of $z=Q_{s+1}(c z^a)$, the values $z$, $c z^a$ and
$Q_{s+2}(c z^a)$ are transcendental.
\end{corollary}
\begin{proof}
If $z$ were algebraic, $t=c z^a$ would be nonzero algebraic and
$e^t=S_s(t)/(1-z)$ would be algebraic. This contradicts
Hermite--Lindemann. Moreover
\[
 Q_{s+2}(c z^a)=1-(1-z)\frac{S_{s+1}(c z^a)}{S_s(c z^a)} .
\]
As a rational function of $z$ its values at $0,1$ are $0,1$, respectively.
The polynomial $c z^a$ is nonconstant as well. Apply
Lemma~\ref{lem:descent}. The conclusion applies to every positive solution
of the fixed-point equation, so it requires no choice between multiple branches.
\end{proof}

\subsection{A different transition in the same hypergraph}
In an $r$-orientation, each hyperedge is assigned to an incident vertex
with at most $r$ assignments per vertex. For $d\ge3$ choices and capacity
$r\ge1$, the threshold in edges per vertex is
\[
 c^*_{d,r}=\frac{T}{d Q_r(T)^{d-1}},\qquad
 \frac{T Q_r(T)}{Q_{r+1}(T)}=dr .
\]
The threshold identification for $r=1$ follows from the cuckoo-hashing and
XORSAT theorem \cite{gCuckoo,gXor}. For $r\ge2$ it is the
multiple-orientability theorem \cite{gFKP}. These theorems specify the
unique positive $T$ in the formula.

\par\hypertarget{cat-G11}{}\label{page:G11}\hypertarget{cat-G12}{}\label{page:G12}\hypertarget{cat-G13}{}\label{page:G13}\hypertarget{cat-G14}{}\label{page:G14}\hypertarget{cat-G15}{}\label{page:G15}\hypertarget{cat-G16}{}\label{page:G16}
\begin{theorem}\label{thm:gorient}
The orientation threshold $c^*_{d,r}$ is transcendental for every
$d\ge3$ and $r\ge1$. For $r=1$ this includes the satisfiability threshold of
random $d$-XORSAT in the ratio of equations to variables.
\end{theorem}
\begin{proof}
Set $K=dr$. Since $Q_{r+1}(T)<Q_r(T)$, we have $0<T<K$.
Rearranging the defining equation yields
\[
   e^T=S_{r-1}(T)+\frac{K T^r}{r!(K-T)}
       =\frac{N(T)}{K-T},\qquad
 N(X)=(K-X)S_{r-1}(X)+KX^r/r! .
\]
This excludes algebraic $T$. It gives $Q_r(T)=KT^r/(r!N(T))$, hence
\[
 c^*_{d,r}=\frac{T}{d}
 \left(\frac{r!N(T)}{KT^r}\right)^{d-1}.
\]
Since $N(0)=K$ and $r(d-1)>1$, this rational function has a pole at
zero. Its denominator at $T$ is nonzero. Lemma~\ref{lem:descent} proves the claim.
\end{proof}

For $d=3,r=1$, the two elementary formulas read
$e^T=(3+2T)/(3-T)$ and $c^*_{3,1}=(3+2T)^2/(27T)$.
This gives the chronological record's $0.917935276658\ldots$ directly.
The smaller edge density $0.818469160761\ldots$ in G08 corresponds to
the birth of the $2$-core. Orientability is lost later.

\begin{center}\small
\begin{tabular}{@{}lrrlc@{}}\toprule
ID & Choices $d$ & Capacity $r$ & Observable & Approximation\\\midrule
G11 & 3&1& XORSAT / cuckoo threshold & 0.917935277\\
G12 & 4&1& XORSAT / cuckoo threshold & 0.976770165\\
G13 & 5&1& XORSAT / cuckoo threshold & 0.992438391\\
G14 & 3&2& Orientability threshold & 1.976402828\\
G15 & 4&2& Orientability threshold & 1.996482968\\
G16 & 3&3& Orientability threshold & 2.991857218\\
\bottomrule\end{tabular}
\end{center}

\subsection{Search trees, extinction, waiting times, and spectral peaks}
\par\hypertarget{cat-G17}{}\label{page:G17}\hypertarget{cat-G18}{}\label{page:G18}
\begin{proposition}[The two scales of a binary search tree]\label{prop:gbst}
G17 and G18 are the height coefficient $4.311070407\ldots$ and the
saturation-level coefficient $0.373364618\ldots$ of a binary search tree
built from a uniform random permutation. Both are transcendental.
\end{proposition}
\begin{proof}
Devroye's theorem \cite{gDev} identifies the two coefficients, measured
relative to $\log n$, as the two positive roots of
$x\log(2e/x)=1$. For such a root
$e^{1/x-1}=2/x$. Since $x=1$ does not solve the equation, an algebraic
$x$ would give a nonzero algebraic exponent and an algebraic exponential.
\end{proof}

\par\hypertarget{cat-G19}{}\label{page:G19}
\begin{proposition}[Poisson branching]\label{prop:gbranch}
For algebraic $c>1$, the survival probability $p(c)$ of a Poisson$(c)$
Galton--Watson process is transcendental. G19 is $p(2)=0.7968121300\ldots$.
\end{proposition}
\begin{proof}
Conditioning on the children of the initial particle gives
$1-p=\exp(-cp)$. Equivalently the extinction probability is the limit
of iterations of $s\mapsto e^{c(s-1)}$ starting at zero, hence its
least fixed point in $[0,1]$. Strict convexity and the derivative $c>1$
at $1$ locate this fixed point in $(0,1)$. Algebraic $p$ would make
both $-cp\ne0$ and $e^{-cp}=1-p$ algebraic.
\end{proof}

\par\hypertarget{cat-G20}{}\label{page:G20}\hypertarget{cat-G21}{}\label{page:G21}\hypertarget{cat-G22}{}\label{page:G22}
\begin{proposition}[Integer-shape gamma quantiles]\label{prop:gerald}
For every integer $h\ge1$ and algebraic $0<p<1$, the $p$-quantile of
a unit-rate gamma distribution of shape $h$ is transcendental.
This includes G20, the shape-$2$ median $1.678346990\ldots$,
G21, the shape-$3$ median $2.674060314\ldots$, and G22, the
shape-$5$ 95th percentile $9.153519027\ldots$.
\end{proposition}
\begin{proof}
Integration by parts in the gamma density $x^{h-1}e^{-x}/(h-1)!$
gives the positive tail $e^{-x}S_{h-1}(x)$. Its strictly negative
derivative on $x>0$ gives a unique positive quantile $t$. Its defining
equation gives $e^t=S_{h-1}(t)/(1-p)$, excluding algebraic $t$.
\end{proof}

\par\hypertarget{cat-G23}{}\label{page:G23}\hypertarget{cat-G24}{}\label{page:G24}\hypertarget{cat-G25}{}\label{page:G25}
\begin{proposition}[Dimensionless Planck peaks]\label{prop:gplanck}
G23, G24 and G25 are the dimensionless peak coordinates of blackbody
energy radiance per unit frequency, wavelength and logarithmic frequency,
respectively. They are $2.821439372\ldots$, $4.965114232\ldots$ and
$3.920690395\ldots$, and each is transcendental.
\end{proposition}
\begin{proof}
Write $x=h\nu/(k_{\rm B}T)=hc/(\lambda k_{\rm B}T)$. For fixed positive
temperature, Planck's radiance formula is proportional to
$\nu^3/(e^x-1)$ per unit frequency. The wavelength Jacobian
$|d\nu/d\lambda|=c/\lambda^2$ gives a factor proportional to $x^5$ in
wavelength coordinates. Multiplying the frequency density by $\nu$ gives
$x^4$ for a logarithmic interval. For $a=3,5,4$, respectively, maximize
$x^a/(e^x-1)$. Its derivative vanishes when $x=a(1-e^{-x})$.
The difference $a(1-e^{-x})-x$ is strictly concave, has positive slope at
zero and is negative at $a$. Thus there is exactly one positive extremum,
a maximum in $(0,a)$. There $e^x=a/(a-x)$ rules out algebraic $x$.
\end{proof}


\section{Quantiles, diffraction, and elliptic geometry}
\label{sec:analytic}
This section gives twenty-five arithmetic conclusions. In particular the
normalization of an optical radius, the branch of an incomplete integral,
and the units of an elliptic period form part of the definition.

\subsection{Distribution quantiles and an exponential-integral zero}
\par\hypertarget{cat-A01}{}\label{page:A01}\hypertarget{cat-A02}{}\label{page:A02}\hypertarget{cat-A03}{}\label{page:A03}\hypertarget{cat-A04}{}\label{page:A04}
\begin{proposition}[A01--A04. Three beta quantiles and a Student quantile]
The $1/2$, $1/40$, and $39/40$ quantiles of the $\operatorname{Beta}(3/2,5/2)$
law (A01--A03) are transcendental. Their approximate values are
$0.3524523$, $0.0387476$, and $0.8232639$, respectively. The $39/40$
quantile of Student's $t$ distribution with three degrees of freedom
(A04), approximately $3.1824463$, is also transcendental.
\end{proposition}
\begin{proof}
For the beta law, direct differentiation and evaluation at zero give its
distribution function $B$ as
\[
 \pi B(x)=2\arcsin\sqrt{x}+\sqrt{x(1-x)}P(x),\qquad
 P(x)=-\tfrac23(4x-1)(2x-3).
\]
If $x$ and $B(x)=r$ were algebraic with $r$ one of the three specified
rationals, $e^{i\pi r}$ would be a root of unity and
$e^{2i\arcsin\sqrt{x}}=(\sqrt{1-x}+i\sqrt{x})^2$ algebraic. Thus
$\exp(i\sqrt{x(1-x)}P(x))$ would be algebraic. Hermite--Lindemann forces
$P(x)=0$ in the open interval. Its sole interior zero is $1/4$ and
$B(1/4)=1/3$, excluding all three probabilities.
For Student's law set $s=x/\sqrt3$. Integration of the normalized density
gives $\pi(F(x)-1/2)=\arctan s+s/(1+s^2)$. For the indicated quantile
$s>0$. Algebraic $s$ would make both $e^{2i\arctan s}=(1+is)/(1-is)$
and $e^{2\pi i(F(x)-1/2)}$ algebraic, hence $e^{2is/(1+s^2)}$ algebraic,
again impossible.
\end{proof}

Here is the specialization input used for the Gaussian and exponential
integrals. Define, for rational $0<a<1$,
\[
 F_a(z)=\sum_{n\ge0}\frac{(-z)^n}{n!(n+a)},\qquad
 E(z)=\operatorname{Ein}(z)=\sum_{n\ge1}\frac{(-1)^{n-1}z^n}{n\,n!}.
\]
\begin{lemma}\label{lem:analytic-multipoint}
For fixed $a$ and distinct nonzero algebraic $b_1,\ldots,b_s$,
the $F_a(b_j)$ are algebraically independent over the field generated
over $\overline{\mathbb Q}$ by the exponentials of algebraic numbers.
The same conclusion holds for the $E(b_j)$.
\end{lemma}
\begin{proof}
It suffices to adjoin finitely many extra exponentials. Take a
$\mathbb Z$-basis $\nu_l$ of the group generated by their exponents and
the $b_j$. Write $a=r/M$ with $0<r<M$. In the rational function field
$K=\mathbb C(t)(e^{\nu_lt^M}:l)$ the exponential generators are
algebraically independent. Indeed, different monomials have distinct
exponents and comparison along a generic ray separates their exponential
growth. The primitives $t^rF_a(b_jt^M)/M$ have derivatives
$t^{r-1}e^{-b_jt^M}$. By the antiderivative theorem of
Ostrowski~\cite{aOst}, dependence would give a nonzero constant linear
combination of these derivatives equal to $R'$ for some $R\in K$.
All fields here are fields of meromorphic germs with constants $\mathbb C$.
Expand $R$ in iterated Laurent series in the exponential generators over
$\mathbb C(t)$. Differentiation preserves characters. At a character
$e^{-b_jt^M}$ with nonzero coefficient $c$ on the right one would have
\[
 A'-Mb_jt^{M-1}A=c t^{r-1},\qquad A\in\mathbb C(t).
\]
A finite pole is impossible by its order in $A'$. A nonzero polynomial
on the left has resulting degree at least $M-1$, whereas $r-1<M-1$.
For $E(b_jt)$ use $M=1$ and $dE(b_jt)/dt=(1-e^{-b_jt})/t$.
At the nontrivial character the analogous equation is $A'-b_jA=-c/t$.
A pole of $A$ at zero would give a pole of order at least two, and a
regular $A$ would give none. This proves functional independence over
the exponential field, after returning to the original variable.
Both displayed series are $E$-functions. Their exponential Taylor
coefficients have exponential size and their common denominators up to
degree $n$ are bounded exponentially by least common multiples up to
$Mn+M$. Include the exponentials of the basis and of all forcing terms
in the differential system. The equations
$zF_a'(z)+aF_a(z)=e^{-z}$ and $zE'(z)=1-e^{-z}$ show that its only
possible finite pole is zero. The Siegel--Shidlovskii theorem at $1$
preserves the total transcendence degree. The forcing exponentials are
Laurent monomials in the basis exponentials, so the asserted relative
independence follows~\cite{aBeu}.
\end{proof}

\par\hypertarget{cat-A05}{}\label{page:A05}\hypertarget{cat-A06}{}\label{page:A06}\hypertarget{cat-A07}{}\label{page:A07}\hypertarget{cat-A08}{}\label{page:A08}
\begin{proposition}[A05--A08. Gaussian medians and entire-function zeros]
The following are transcendental.
\begin{enumerate}
\item[\textbf{A05.}] The median on $[-1,1]$ of the normalized density
proportional to $e^{-(x+1/2)^2}+2e^{-(x-1/2)^2}$, about $0.1111020$.
\item[\textbf{A06.}] The median of $e^{-x^2}$ normalized on $[0,1]$,
about $0.3926976$.
\item[\textbf{A07.}] The unique $t>1$ solving
$E(t^2)-E(t)=t(t-1)/2$, about $1.3569857$.
\item[\textbf{A08.}] The first upper-half-plane zero of
$\operatorname{erfc}$ tabulated in \cite[\S7.13]{aDLMF},
about $-1.3548101+1.9914668i$. In fact every zero is transcendental.
\end{enumerate}
\end{proposition}
\begin{proof}
Let $G(x)=\int_0^xe^{-s^2}\,ds=xF_{1/2}(x^2)/2$. At distinct positive
algebraic radii its values are linearly independent over
$\overline{\mathbb Q}$ by Lemma~\ref{lem:analytic-multipoint}.
The first median satisfies
\[
 G(x+1/2)+2G(x-1/2)=\tfrac12G(1/2)-\tfrac12G(3/2).
\]
Neither radius on the left reaches $3/2$ in the open interval.
After applying oddness, the nonzero coefficient of $G(3/2)$ therefore
excludes algebraic $x$. The second median satisfies $G(x)=G(1)/2$ with
$0<x<1$, excluded by the same independence.
For A07 the quotient $(E(t^2)-E(t))/(t^2-t)$ is the mean of
$E'(t+s(t^2-t))$ for $s\in[0,1]$. Since
$E'(u)=\int_0^1 e^{-vu}\,dv$, this quotient decreases strictly from
$1-e^{-1}$ to zero on $(1,\infty)$. The unique root is transcendental by
the lemma at its two distinct arguments, under the algebraicity assumption.
For A08 real zeros are excluded by positivity of the Gaussian tail,
and purely imaginary zeros by $\Re\operatorname{erfc}(iy)=1$.
At any remaining hypothetical algebraic zero $z$, the numbers
$z^2$ and $\bar z^2$ are different and nonzero. Yet
$zF_{1/2}(z^2)=\bar zF_{1/2}(\bar z^2)=\sqrt\pi$. This contradicts
the lemma.
\end{proof}

The same argument applies on $\mathbb C\setminus(-\infty,0]$ to
every zero of the upper incomplete gamma function $\Gamma(a,z)$
for any rational nonintegral $a$. For $0<a<1$ a zero would satisfy
$z^aF_a(z)=\Gamma(a)$, a real constant. Conjugation then gives a
forbidden relation at $z,\bar z$ if $z$ is nonreal algebraic.
At positive real $z$ the upper-tail integral is positive even for
negative $a$. For other nonintegral rational shapes the recurrence
$\Gamma(a+1,z)=a\Gamma(a,z)+z^ae^{-z}$ transfers a zero to
a nonzero affine relation in the two independent $F_{\{a\}}$ values,
with coefficients in the exponential field. This family is unnumbered.

\subsection{Bessel extrema, exponential fixed points, and optical radii}
\par\hypertarget{cat-A09}{}\label{page:A09}\hypertarget{cat-A10}{}\label{page:A10}
\begin{proposition}[A09--A10. Modified Bessel extremum]
For $R(y)=I_1(y)/I_0(y)$ put $D(y)=(yR(y))'$. There is a unique
positive maximum point (A09) and its value (A10) is greater than one.
Both numbers are transcendental. More generally the maximum point and
value have these arithmetic properties when
$R=I_{\mu+1}/I_\mu$ for rational $-1<\mu<1/2$.
\end{proposition}
\begin{proof}
Write $a=\mu+1/2$, $v=R+a/y$. Bessel's recurrence yields
\[
 v'=1-v^2+a(a-1)/y^2=(D-v)/y,\qquad D'=2(1-vD).
\]
Positivity $D>0$ follows by expanding $H=y^{-\mu}I_\mu(y)$ up to a
constant as a series in $y^{2n}$ with positive coefficients. The weights
proportional to these terms give $D=4\operatorname{Var}(n)/y$.
Near zero $D\sim y/(\mu+1)$ and at infinity
$D=1-a(a-1)/(2y^2)+O(y^{-3})$, with differentiable asymptotics
\cite[\S10.40]{aDLMF}. At a zero of $v'$,
$v''=-2a(a-1)/y^3$. At a zero of $D'$ one has $v=1/D$ and
$D''=-2(D^2-1)/y$. The value $D=1$ there would force $a(a-1)=0$.
For $0<a<1$, $v'$ starts negative and ends positive. Each of its
zeros is an upward crossing. Before the unique crossing $D<v$, so a
stationary point has $D<1$ and would be a minimum, incompatible with
a first stationary point since $D'$ starts positive. Afterwards every
stationary point is a maximum with $D>1$, and $D'$ eventually becomes
negative. For $a<0$ the argument reverses for $v'$. It starts positive,
ends negative, and crosses once. Before the crossing only maxima above
one, and after it only minima below one, are possible. An overshoot
$D>1$ must occur since $y(1-R)$ starts at zero, has derivative $1-D$,
and tends to $a<0$. The far tail approaches one from below, giving a
unique maximum and a later minimum. For $a=0$, $R=\tanh y$ and
$D'=2\operatorname{sech}^2 y(1-y\tanh y)$. At its sole positive stationary point $D=y>1$.

At a stationary location $y$ the relation
\[
 y-(yR+a)\{y(1-R^2)-2\mu R\}=0
\]
is cubic in $R$ with leading coefficient $y^2$. For nonzero algebraic
$y$, the normalized function $H(y)=\Gamma(\mu+1)(y/2)^{-\mu}I_\mu(y)$ and its derivative are
algebraically independent by Siegel's Bessel theorem
\cite{aSiegel,aFeld}. Here $R=H'/H$. This applies for
$\mu\notin\tfrac12+\mathbb Z$. The remaining $\mu=-1/2$ gives
$R=\tanh y$, transcendental at positive algebraic $y$ by
Hermite--Lindemann. Thus the location is transcendental. Eliminating
$v=1/m$ at the maximum, where $m=D(y)>1$, gives
\[
 (m^2-1)y^2-m(m^2-1)y+(\mu^2-1/4)m^2=0.
\]
Algebraic $m$ would contradict the transcendence of $y$.
\end{proof}

\par\hypertarget{cat-A11}{}\label{page:A11}\hypertarget{cat-A12}{}\label{page:A12}\hypertarget{cat-A13}{}\label{page:A13}\hypertarget{cat-A14}{}\label{page:A14}
\begin{proposition}[A11--A14. Laplace, Bratu, Steklov, and Dawson]
Let $u>0$ solve $u\tanh u=1$. The Laplace limit of Kepler's equation
$L=1/\sinh u$ (A11) and the one-dimensional unit-interval Bratu critical
parameter $b=8(u^2-1)$ (A12) are transcendental. The optimal normalized
first Steklov eigenvalue of an annulus $S=4\pi/u$ (A13) has irrational
square. If $D_0(x)=e^{-x^2}\int_0^xe^{s^2}\,ds$ is Dawson's integral,
$c=2\max_{x\ge0}D_0(x)$ (A14) is transcendental.
\end{proposition}
\begin{proof}
The function $u\tanh u$ increases from zero to infinity. Its root has
$u>1$ and $e^{2u}=(u+1)/(u-1)$. Hermite--Lindemann forces $u$ to be
transcendental. The defining formulas in \cite{aMcM,aBuck,aFraser}
give $L^2=u^2-1=b/8$, proving both transcendence statements. If
$S^2$ were rational then $u/\pi=\sqrt{d}/h$ for positive integers
$d,h$. The relation for $e^{2u}$ would make $e^{\pi\sqrt d}$ algebraic
over $\overline{\mathbb Q}(\pi)$, contradicting Nesterenko's
independence of $\pi,e^{\pi\sqrt d}$ \cite[Corollary 6]{aNest}.
For Dawson, differentiation gives $D_0'=1-2xD_0$. Positivity and
$D_0(x)\to0$ as $x\to\infty$ give an interior maximum $t$. Then
$D_0''(t)=-2D_0(t)<0$ at every stationary point implies uniqueness.
The identity $D_0(x)=e^{-x^2}xF_{1/2}(-x^2)/2$ and
Lemma~\ref{lem:analytic-multipoint} imply $D_0(x)$ transcendental
at nonzero algebraic real $x$. Hence $D_0(t)=1/(2t)$ excludes algebraic
$t$, and $c=1/t$ is transcendental.
\end{proof}
\par\hypertarget{cat-A15}{}\label{page:A15}\hypertarget{cat-A16}{}\label{page:A16}\hypertarget{cat-A17}{}\label{page:A17}
\begin{proposition}[A15--A17. Circular and annular diffraction]
In the dimensionless Bessel coordinate, the Airy half-energy radius
(A15, about $1.6802247$), the half-energy radius of a pupil with
central obscuration ratio $1/2$ (A16, about $3.9371301$), and the first
dark-ring radius of the unobscured pupil (A17, $j_{1,1}\approx3.8317060$)
are transcendental.
\end{proposition}
\begin{proof}
For $0\le h<1$ the focal-plane encircled-energy fraction is
\[
 B_h(x)=\frac{2}{1-h^2}\int_0^x
       \frac{(J_1(t)-hJ_1(ht))^2}{t}\,dt .
\]
The usual radial Fourier transform of the annular pupil gives this
formula \cite{aPupil}. The normalization at infinity also follows from Hankel
Plancherel applied to $\mathbf1_{(h,1)}(r)$, whose order-zero transform
is $(J_1(t)-hJ_1(ht))/t$. The integrand has only isolated zeros on
$(0,\infty)$, giving a unique radius at fraction $1/2$. Differentiating
$J_0^2+J_1^2$ gives $-2J_1^2/x$. Consequently
$B_0=1-J_0^2-J_1^2$. At a nonzero algebraic argument $J_0,J_1$
are algebraically independent by Siegel's Bessel theorem. This rules
out both $B_0(x)=1/2$ and $J_1(x)=0$, proving A15 and A17.

Here is the extra argument needed for the annular pupil. Put $h=1/2$,
$A_j=J_j(x)$, $B_j=J_j(hx)$ for $j=0,1$, and
$C(x)=\int_0^x A_1(t)B_1(t)\,dt/t$ (where $A_1(t)=J_1(t)$ and
$B_1(t)=J_1(ht)$). The exact equation is
\[
 (1-h^2)B_h(x)=1+h^2-A_0^2-A_1^2-h^2(B_0^2+B_1^2)-4hC(x).
\]
Consider the vector formed by $1$, the ten quadratic monomials in
$A_0,A_1,B_0,B_1$, and $C$. These are $E$-functions satisfying a
first-order linear system with possible finite pole only at zero.
For $C$, the product has zero constant term and integration against
$dt/t$ divides each nonconstant exponential Taylor coefficient by its
index, preserving the exponential common-denominator bound.
The four Bessel functions are algebraically independent over
$\overline{\mathbb Q}(x)$. A polynomial relation could be specialized
at a nonzero rational point avoiding coefficient zeros, contrary to
Siegel's simultaneous Bessel theorem for arguments $x,hx$ of different
squares \cite[\S4, Theorem 2]{aFeld}.

Let $W$ be the rational-function span of $1$ and these monomials. Any
linear relation among the extended vector has zero coefficient on
$A_0^2$. Indeed, if $C\notin W$ this follows by independence. If
$C\in W$, its unique representation has a part $P_A$ involving only
the three pure $A$ monomials. Differentiation preserves the four
summands (constant, pure $A$, mixed, pure $B$) by Bessel's equations.
As $C'$ is mixed, $P_A'=0$. Thus $P_A$ is constant, which forces
$P_A=0$ by algebraic independence. Every functional relation is then
a multiple of this representation and has zero $A_0^2$ coefficient.
If $B_h(\alpha)=1/2$ at a nonzero algebraic $\alpha$, the energy
equation would give an algebraic linear relation with $A_0^2$
coefficient $-1$ there. Beukers's lifting theorem at the ordinary point
$\alpha$ lifts the relation with the same specialized coefficients
\cite[Theorem 1.3]{aBeu}. This is impossible.
\end{proof}
The arithmetic assertions concern the coordinate $x$. Expressing
a radius in the optical unit $\lambda/D$ introduces the factor $1/\pi$.

\subsection{Square periods and rotational minimal hypersurfaces}
\par\hypertarget{cat-A18}{}\label{page:A18}
\begin{proposition}[A18. A square-lattice sum]
With $z=j+ik$ the absolutely convergent sum
\[
 C=\sum_{j,k\in\mathbb Z}
   \frac{1}{(4z+1)(4z-1)(4z-3)}
   =0.3956379318\ldots
\]
is transcendental and satisfies
$C=(1+\sqrt2)T/16$, where $T$ is the least positive real period of the
Weierstrass function with invariants $(4,0)$.
\end{proposition}
\begin{proof}
The period lattice is $T(\mathbb Z+i\mathbb Z)$, with half-period
$h=T/2$, $\wp(h)=1$, and $\wp'(h)=0$. Put $a=T/4$. Writing $\omega=Tz$ puts the middle denominator at $(\omega-a)/a$.
Reflecting $\omega$ in the lattice changes the centered triple reciprocal
to $-1/[w(w-h)(w+h)]$, where $w=\omega+a$, with the scale factor $a^3$.
Partial fractions and the addition formula give
\[
 C=-\frac{a^3}{2h^2}\{\zeta(a+h)+\zeta(a-h)-2\zeta(a)\}
   =-\frac{a^3}{2h^2}\frac{\wp'(a)}{\wp(a)-1}.
\]
Here $\zeta$ is Weierstrass's zeta function. The convergence corrections
in its partial-fraction expansion cancel in the second difference.
The half-period shift formula and $\wp(a+h)=\wp(a)$ give
$(\wp(a)-1)^2=2$. On the positive real quarter-period $\wp(a)>1$ and
$\wp'(a)<0$, so $\wp(a)=1+\sqrt2$ and
$\wp'(a)=-2\sqrt2(1+\sqrt2)$. Substitution gives the stated multiple
of $T$. A nonzero elliptic period for algebraic invariants is
transcendental by Schneider's theorem \cite{aSch,aWald}.
\end{proof}

\par\hypertarget{cat-A19}{}\label{page:A19}
\begin{proposition}[A19. A four-dimensional catenoid]
The maximal gap divided by boundary-sphere radius for a rotational
three-dimensional catenoid between equal coaxial spheres in
$\mathbb R^4$ is transcendental. Its value is
$\Lambda_3\approx0.8447865286$.
\end{proposition}
\begin{proof}
Scaling the neck radius to one, the rotational area integral has
Lagrangian $r^2\sqrt{1+(r')^2}$ in the axial coordinate. Its first
integral is $r^2/\sqrt{1+(r')^2}=1$. The gap/radius at boundary radius
$u>1$ is therefore $2I(u)/u$, with
$I(u)=\int_1^u ds/\sqrt{s^4-1}$. It tends to zero at either end and
has one maximum. Indeed $uI'(u)-I(u)$ decreases strictly from positive
infinity to a negative limit because $I''<0$. At the maximizer,
\[
 I(u)=u/\sqrt{u^4-1},\qquad \Lambda_3=2/\sqrt{u^4-1}.
\]
If $u$ were algebraic, $I(u)>0$ would be a nonzero elliptic integral
of the first kind with algebraic endpoints on $Y^2=X^4-1$. It would
be transcendental by Schneider's theorem, contrary to its displayed
value. Algebraic $\Lambda_3$ would imply algebraic $u$, completing
the proof.
\end{proof}

\subsection{Elastic curves and a length ratio}
Let $F(\phi,m),E_2(\phi,m)$ denote the incomplete elliptic integrals
of the first and second kinds, respectively, with parameter $m$
(the square of the modulus). Write $K(m)=F(\pi/2,m)$ and
$\alpha(m)=\arcsin(1/\sqrt{2m})$ for $1/2<m<1$.
\par\hypertarget{cat-A20}{}\label{page:A20}\hypertarget{cat-A21}{}\label{page:A21}\hypertarget{cat-A22}{}\label{page:A22}
\begin{proposition}[A20--A22. Two elastica parameters and a length ratio]
The teardrop and heart parameters $m_T\approx0.7311830$ (A20)
and $m_H\approx0.8436273$ (A21) are transcendental. The teardrop
arc length divided by the half-period of its generating elastic wave
(A22) is the irrational number
\[
 r_T=\frac{F(\pi-\alpha(m_T),m_T)}{K(m_T)}\approx1.4822533 .
\]
\end{proposition}
\begin{proof}
Miura, M\"uller and Rupp \cite[\S2, Definitions 2.15 and 2.21]{aMMR}
prove existence and uniqueness in the geometric ranges of these parameters.
For the teardrop put $\phi=\pi-\alpha(m_T)$ and $u=F(\phi,m_T)$.
Its defining equation is $2E_2(\phi,m_T)-u=0$. The symmetric arc has
length $2u$ in the arclength variable and the generating wave has
half-period $2K$. For the heart the defining equation is
\[
 \int_{-\pi/4}^{5\pi/4}
       \frac{1-2\sin^2\theta}{\sqrt{1-m_H\sin^2\theta}}\,d\theta=0 .
\]

Suppose in either case that $m$ is algebraic. Use Jacobi functions
with parameter $m$, set $c=(2-m)/3$, and choose Weierstrass functions
$\wp(v)=1/\operatorname{sn}^2v-(1+m)/3$, $\zeta'=-\wp$. Their lattice invariants are algebraic functions of $m$ with rational coefficients. Put $\epsilon(v)=E_2(\operatorname{am}v,m)$.
Differentiation and the common normalization at zero give
\[
 \epsilon(v)=\zeta(v)+cv-
 \frac{\operatorname{cn}v\operatorname{dn}v}{\operatorname{sn}v}.
\]
At the teardrop endpoint the Jacobi functions are algebraic and
$\operatorname{sn}u\ne0$. Thus $\wp(u)$ and $(2c-1)u+2\zeta(u)$
would both be algebraic. Schneider's elliptic value theorem
\cite[Theorem 20(1)]{aWald} excludes simultaneous algebraicity of
$\wp(v)$ and $av+b\zeta(v)$ for algebraic $(a,b)\ne(0,0)$
away from lattice points.

For the heart take $v=F(\pi/4,m)$ and
$w=2K+2v$. The arclength parameters of the two integration endpoints
are $-v$ and $2K+v$. Multiplying the heart equation by $m$ gives
$2[\epsilon(2K+v)-\epsilon(-v)]+(m-2)w=0$.
Their Jacobi functions are algebraic. Also $0<v<K$, so $w$ is
strictly between $2K$ and $4K$ and is not a pole. Duplication gives
algebraic $\wp(w)=\wp(2v)$, and
\[
 \zeta(2K+v)-\zeta(-v)-\zeta(w)
     =2\zeta(v)-\zeta(2v)=-\wp''(v)/(2\wp'(v))
\]
is algebraic, since $\wp'(v)\ne0$. It follows that
$(2c+m-2)w+2\zeta(w)$ would be algebraic, again impossible.

To prove A22, uniformize the Legendre curve
$y^2=x(x-1)(x-m)$ by $x=1/\operatorname{sn}^2u$ and
$y=x'(u)/2$. At the teardrop endpoint $x=2m$ and
$y^2=2m^2(2m-1)$. If $u/K$ were rational, this point would be torsion,
since $2K$ is a period. Fix a putative torsion order $N$. The
$N$-division polynomial evaluated at $x=2m$ gives a polynomial in $m$
with rational coefficients (for even $N$ first remove the factor
$2y$, nonzero here). This polynomial is nonzero. To see this, specialize to
$m=3/2$, giving $x=3$, $y=\pm3$. On the integral short Weierstrass model
$X=36x-30$, $Y=216y$, one obtains
$Y^2=X^3-756X+4320$, with $P=(78,648)$ up to sign.
The tangent formula gives $2P=(105/4,405/8)$, whose coordinates are
not integral. The Nagell--Lutz theorem \cite[Chap.~VIII, \S7]{aSil}
excludes rational torsion for $2P$, hence for $P$. Consequently a torsion
relation at $m_T$ would force the transcendental number $m_T$ to be a root
of a nonzero rational polynomial, a contradiction.
\end{proof}

\subsection{Oscillatory spectra}
\par\hypertarget{cat-A23}{}\label{page:A23}\hypertarget{cat-A24}{}\label{page:A24}\hypertarget{cat-A25}{}\label{page:A25}
\begin{proposition}[A23--A25. Beam roots and a diffraction sidelobe]
The first positive solution of $\cosh x\cos x=-1$ (A23), the first
positive solution of $\cosh x\cos x=1$ (A24), and the unique solution
of $\tan x=x$ in $(\pi,3\pi/2)$ (A25) are transcendental.
The first two are the dimensionless clamped--free and clamped--clamped
Euler--Bernoulli beam roots. The third gives the position of the first
positive sidelobe maximum of $(\sin x/x)^2$ after its first zero.
\end{proposition}
\begin{proof}
For either beam sign, a positive solution exists by continuity (use
$\pi/2,\pi$ for the negative sign and $3\pi/2,2\pi$ for the positive
sign). Zeros are isolated, including the root at zero for the positive
sign. If $x\ne0$ were algebraic, the four exponentials with exponents
$(\pm1\pm i)x$, together with $e^0$, would be linearly independent
over the algebraic numbers by Lindemann--Weierstrass. Their sum
$4\cosh x\cos x$ therefore cannot equal $\pm4$.
For A25, $\tan x-x$ increases from a negative value to positive infinity
on the stated interval. At its zero $e^{2ix}=(1+ix)/(1-ix)$, excluding
algebraic $x$ by Hermite--Lindemann.
\end{proof}


\section{Modular values, lattice walks and growth rates}\label{sec:mcat}

\subsection{The peak of the modular discriminant}

Write $q=e^{2\pi i\tau}$, $\operatorname{Im}\tau>0$, and $\theta=q\,d/dq$.
Let $E_2,E_4,E_6$ be the Eisenstein series with constant term $1$ and
$\Delta(q)=q\prod_{n\ge1}(1-q^n)^{24}$.
Nesterenko's theorem \cite{aNest} gives
\[
 \operatorname{trdeg}_{\mathbb Q}\mathbb Q(q,E_2(q),E_4(q),E_6(q))\ge3
 \qquad(0<|q|<1).
\]

\par\hypertarget{cat-M01}{}\label{page:M01}\hypertarget{cat-M02}{}\label{page:M02}\hypertarget{cat-M03}{}\label{page:M03}
\begin{proposition}[M01--M03]
Let $q_0$ maximize $\Delta$ on $(0,1)$, let $D=\Delta(q_0)$ and let
$K=-\Delta''(q_0)$, where the derivatives are taken with respect to $q$.
The three numbers $q_0$ (M01), $D$ (M02) and $K$ (M03) are algebraically independent.
\end{proposition}
\begin{proof}
The logarithmic derivative $\theta\Delta/\Delta=E_2
=1-24\sum_{n\ge1}nq^n/(1-q^n)$ decreases strictly from $1$ to $-\infty$.
Its unique zero is the maximum. There $q_0,E_4,E_6$ are algebraically
independent by Nesterenko. The identities
$\theta E_2=(E_2^2-E_4)/12$ and $1728\Delta=E_4^3-E_6^2$ give
\[
 E_4(q_0)=12q_0^2K/D,\qquad E_6(q_0)^2=E_4(q_0)^3-1728D.
\]
Thus $\mathbb Q(q_0,D,K)$ has transcendence degree at least three.
\end{proof}

\subsection{Cubic-lattice walks and electrical resistance}

Let a walk choose each nearest neighbour with equal probability. Denote its
probability of a positive-time return to the origin by $p$. If $W$ is the expected
number of visits, counting time zero, the renewal identity is $W=1+pW$.

\par\hypertarget{cat-M04}{}\label{page:M04}\hypertarget{cat-M05}{}\label{page:M05}\hypertarget{cat-M06}{}\label{page:M06}\hypertarget{cat-M07}{}\label{page:M07}\hypertarget{cat-M08}{}\label{page:M08}
\begin{proposition}[M04--M08]
The simple cubic return probability $p_{\mathrm{sc}}$ (M04), the simple cubic unit-edge
effective resistances $A=R(0,(1,1,0))$ (M05) and $B=R(0,(1,1,1))$ (M06),
and the face-centred and body-centred cubic return probabilities $p_{\mathrm{fcc}}$
(M07) and $p_{\mathrm{bcc}}$ (M08) are transcendental.
Moreover, $A,B$ are algebraically independent, as are
$p_{\mathrm{sc}},\pi,e^{\pi\sqrt6}$.
\end{proposition}
\begin{proof}
For the simple cubic lattice, Fourier inversion gives $W=3g$, where
\[
 g=\pi^{-3}\int_{[0,\pi]^3}\frac{du\,dv\,dw}{3-\cos u-\cos v-\cos w}.
\]
Watson's evaluation, in the normalization of \cite{mGlasser}, reads
\[
 g=4(18+12\sqrt2-10\sqrt3-7\sqrt6)\,K(k)^2/\pi^2,
 \qquad k=(2-\sqrt3)(\sqrt3-\sqrt2).
\]
Here $K(k)=\int_0^{\pi/2}(1-k^2\sin^2t)^{-1/2}dt$, and $k$ is the
singular modulus with $K(\sqrt{1-k^2})/K(k)=\sqrt6$.
The corresponding elliptic curve has algebraic invariants and complex
multiplication by $\mathbb Q(\sqrt{-6})$. The integral $K(k)$ is an algebraic multiple
of a period. Nesterenko's CM-period corollary \cite{aNest} gives algebraic
independence of $\pi,e^{\pi\sqrt6},K(k)$.
The formula for $g$ then gives the same independence with $K(k)$ replaced by
$g$, and renewal gives $p_{\mathrm{sc}}=1-1/(3g)$.

For a unit resistor network, Fourier inversion of its Laplacian gives
$R(0,v)=g-G(3;v)$ with the same Green function normalization.
The evaluations of $G(3;1,1,0)$ and $G(3;1,1,1)$ in \cite{mGlasser} give
\[
 A=\tfrac7{12}g+\tfrac12 h,\qquad
 B=\tfrac98g-\tfrac34 h,\qquad h=(\pi^2g)^{-1}.
\]
The two-by-two coefficient matrix has determinant $-1$. Since $\pi,g$
are algebraically independent, so are $g,h$ and therefore $A,B$.
In fact $g=3A/4+B/2$, so $p_{\mathrm{sc}},A,B$ satisfy a relation.

For the other two lattices, use the walk normalization of \cite{mGutt} and the elliptic evaluations in \cite[eqs.~(17), (27)]{mGlasser}. These give
\[
 W_{\mathrm{fcc}}=\frac{9\,\Gamma(1/3)^6}{2^{14/3}\pi^4},\qquad
 W_{\mathrm{bcc}}=\frac{\Gamma(1/4)^4}{4\pi^3}.
\]
Each of the pairs $(\pi,\Gamma(1/3))$ and $(\pi,\Gamma(1/4))$ is
algebraically independent \cite{aNest}. Each $W$ is thus transcendental,
and $p=1-1/W$ finishes both proofs.
\end{proof}

\subsection{Growth constants from three types of generating function}

For a set $S$ of positive integers, a composition is an ordered sequence of
parts in $S$. Write $A(z)=\sum_{s\in S}z^s$. The composition generating
function, including the empty sequence, is $(1-A(z))^{-1}$. More generally,
$A$ can count a class of possible blocks with multiplicities.
In each application below $A$ has nonnegative coefficients, is analytic
on $|z|<1$, has $A(0)=0$, tends to $+\infty$ as $z\uparrow1$, and has support
of gcd one. The unique $\rho\in(0,1)$ satisfying $A(\rho)=1$ is a simple
pole of $(1-A)^{-1}$. Equality in the triangle inequality shows that another
pole with $|z|=\rho$ would require $(z/\rho)^s=1$ for every supported $s$.
The gcd condition excludes this. Isolated zeros and compactness give
$r\in(\rho,1)$ on whose closed disk there is no further pole. Subtracting
the principal part and applying Cauchy's estimate yields
\[
 [z^n](1-A(z))^{-1}=\frac{\rho^{-n}}{\rho A'(\rho)}+O(r^{-n}).
\]
In particular, the exponential growth constant is $\rho^{-1}$.

\par\hypertarget{cat-M09}{}\label{page:M09}\hypertarget{cat-M10}{}\label{page:M10}\hypertarget{cat-M11}{}\label{page:M11}\hypertarget{cat-M12}{}\label{page:M12}
\begin{proposition}[M09--M12]
For irrational $a>1$ let $\lambda_a$ be the composition growth constant with parts
$\lfloor ka\rfloor$, $k\ge1$. For $\varphi=(1+\sqrt5)/2$, the four constants
$\lambda_{\varphi^2}$ (upper Wythoff, M09), $\lambda_\varphi$ (lower Wythoff, M10),
$\lambda_{1+\sqrt2}$ (M11) and $\lambda_\pi$ (M12) are transcendental.
\end{proposition}
\begin{proof}
Consecutive parts have differences $\lfloor a\rfloor$ and $\lceil a\rceil$,
and both occur because the average difference tends to the irrational number $a$.
Thus the gcd is one. With $\beta=1/a$ and $G(z)=\sum_{n\ge1}\lfloor n\beta\rfloor z^n$,
the coefficient at $z^n$ in $A$ is $\lfloor(n+1)\beta\rfloor-\lfloor n\beta\rfloor$.
Hence $A(z)=(1-z)G(z)/z$. At its root $\rho$, $G(\rho)=\rho/(1-\rho)$.
The Hecke--Mahler theorem for irrational slope \cite{mBeatty} makes
$G(\rho)$ transcendental at every nonzero algebraic $\rho$ in the unit disk.
This contradiction proves the assertion for each slope. Distinct slopes here
give distinct constants, since every part is nondecreasing with $a$ and
some part strictly increases.
\end{proof}

\par\hypertarget{cat-M13}{}\label{page:M13}\hypertarget{cat-M14}{}\label{page:M14}\hypertarget{cat-M15}{}\label{page:M15}
\begin{proposition}[M13--M15]
Take binary words consisting of any number of pairs of runs, first a zero-run
of length $2^j$, then a one-run of length $b^k$, with $j,k\ge0$.
Their exponential growth constants for $b=3$ (M13), $5$ (M14) and $7$
(M15) are transcendental.
\end{proposition}
\begin{proof}
Put $F_d(z)=\sum_{j\ge0}z^{d^j}$. The generating function of a pair is
$A=F_2F_b$. Its support contains $2$ and $3$.
The functional equation $F_d(z)=z+F_d(z^d)$ gives a Mahler system for
$(1,F_d)$ with polynomial matrix and determinant one. The power series
has the unit circle as a natural boundary by the Hadamard gap theorem.
Nishioka's specialization theorem, as in \cite{mMahler}, thus makes
$F_d(\alpha)$ transcendental at each nonzero algebraic $\alpha$ in the disk, since every such point is regular for this system. By the independent-base
theorem of Adamczewski and Faverjon \cite[Theorem~1.1(ii)]{mMahler},
$F_2(\alpha)$ and $F_b(\alpha)$ are algebraically independent, because
$2$ and $b$ are multiplicatively independent. At $\rho$ their product is
one. This forces $\rho$ and its inverse to be transcendental.
Increasing $b$ strictly decreases $F_b(z)$ for $0<z<1$, so the three
growth constants are distinct.
\end{proof}

\par\hypertarget{cat-M16}{}\label{page:M16}\hypertarget{cat-M17}{}\label{page:M17}\hypertarget{cat-M18}{}\label{page:M18}
\begin{proposition}[M16--M18]
The composition growth constants for the positive generalized pentagonal
numbers $k(3k-1)/2$, $k\in\mathbb Z\setminus\{0\}$ (M16),
positive squares (M17), and positive triangular numbers (M18) are transcendental.
\end{proposition}
\begin{proof}
Each part set contains $1$. Put
\[
 P=\sum_{k\in\mathbb Z}q^{k(3k-1)/2},\qquad
 S=\sum_{k\in\mathbb Z}q^{k^2},\qquad
 U=\sum_{k\ge0}q^{k(k+1)/2}.
\]
The respective root equations are $P(\rho)=2$, $S(\rho)=3$ and $U(\rho)=2$.
With $\Delta_j=\Delta(q^j)$, the Jacobi triple product gives
\[
 qP^{24}=\frac{\Delta_2\Delta_3^2}{\Delta_1\Delta_6},\qquad
 S^{24}=\frac{\Delta_2^5}{\Delta_1^2\Delta_4^2},\qquad
 q^3U^{24}=\frac{\Delta_2^2}{\Delta_1}.
\]
For instance the product for $P$ is
$\prod_{m\ge1}(1-q^{3m})(1+q^{3m-1})(1+q^{3m-2})$.
The three displayed quotients are nonconstant meromorphic modular forms of
weight $12$ on congruence subgroups and have rational Fourier coefficients.
The algebraic-nome theorem of Bhowmik and Pathak
\cite[Theorem~1.8(a)]{mBhowmik} makes their finite nonzero values
transcendental at algebraic nomes in the punctured disk. In each case the
left side at an algebraic $\rho$ would be algebraic, a contradiction.
\end{proof}


The leading coefficients in the preceding proposition also admit a stronger
conclusion. For each of its three part sets, put $C=1/(\rho A'(\rho))$.
Then $C$ and $\rho^{-1}$ are algebraically independent.

\begin{proof}
In each case the displayed modular quotient is $H(q)=q^uV(q)^{24}$, with
$V=1+bA$ and $(u,b)=(1,1),(0,2),(3,1)$, respectively. At the root,
\[
 H(\rho)=\rho^u(1+b)^{24},\qquad
 (\theta H/H)(\rho)=u+\frac{24b}{(1+b)C}.
\]
Write $j=E_4^3/\Delta$. The ratios $E_4^3/H$, $E_6^2/H$ and
$\big((\theta H/H)-E_2\big)^6/H$ have weight zero and algebraic
Fourier coefficients. Indeed $\theta H/H-E_2$ is modular of weight two
by the weight-twelve transformation law. These ratios are finite at the root
and algebraic over $\overline{\mathbb Q}(j(\rho))$ by modular-function-field
theory \cite[Theorem~2.11]{mBhowmik}. Consequently
$\rho,E_2,E_4,E_6$ are algebraic over $\overline{\mathbb Q}(\rho,C,j(\rho))$
when the forms are evaluated at the root. Nesterenko's lower bound of three
forces $\rho,C$ to be algebraically independent.
\end{proof}

\subsection{A level-five zero of reflected \(q\)-Stieltjes constants}

Use the $q$-Hurwitz normalization of \cite[formula~(7)]{nCG}.
For $Q>1$ put $t=Q^{-1/5}$, and for $a=1,2$ put $\eta_a=1-2a/5$. The
normalized reflected differences have the absolutely convergent expressions
\[
 Y_a(Q)=
 \frac{\gamma_0(Q,a/5)-\gamma_0(Q,1-a/5)}{Q-1}
 =\eta_a+\sum_{n\ge1}\frac{t^{an}-t^{(5-a)n}}{1-t^{5n}} .
\]
Only the convergent Laurent-coefficient formula is used here. Set
$H(Q)=Y_1(Q)-4Y_2(Q)$.

\par\hypertarget{cat-M19}{}\label{page:M19}\hypertarget{cat-M20}{}\label{page:M20}\hypertarget{cat-M21}{}\label{page:M21}
\begin{proposition}[M19--M21]
There is exactly one $Q_*$ such that $H(Q_*)=0$ and
$617/1000<Q_*^{-1/5}<618/1000$.
The quantities $Q_*$ (M19), $Y_2(Q_*)$ (M20) and $H'(Q_*)$ (M21) are
algebraically independent. Approximately they are $11.1304720$, $0.5236094$
and $-0.00372877$, respectively.
\end{proposition}
\begin{proof}
For explicit root isolation let $l=617/1000$, $u=618/1000$ and
$c_a(d)=\mathbf{1}_{d\equiv a\ (5)}-\mathbf{1}_{d\equiv-a\ (5)}$. Write
$c=c_1-4c_2$ and $h_n=\sum_{d\mid n}c(d)$. Expanding the series gives
$H=-1/5+\sum_{n\ge1}h_nt^n$. For reproducibility, take
$P(t)=-1/5+\sum_{n=1}^{60}h_nt^n$. Direct rational sums give
\[
 -0.000197<P(l)<-0.000196,\qquad
 0.000139<P(u)<0.000140 .
\]
A lower bound for $P'$ on this interval is
$\sum_{n=1}^{60}n h_n x_n^{n-1}>0.284$, where $x_n=l$ for $h_n\ge0$
and $x_n=u$ otherwise. Since $|h_n|\le4n$, the value and derivative
tails are bounded by $4\sum_{n>60}nu^n<2\cdot10^{-10}$ and
$4\sum_{n>60}n^2u^{n-1}<2\cdot10^{-8}$, respectively.
The change of variable $t=Q^{-1/5}$ is one-to-one.

Now put $t=e^{2\pi i\tau}$ and $\theta=t\,d/dt$. Each
$f_a=Y_a-\eta_a/2$ is a weight-one modular form on $\Gamma_1(5)$.
Indeed the odd function $c_a$ on $(\mathbb Z/5\mathbb Z)^\times$ is a
linear combination of odd Dirichlet characters. Their weight-one Eisenstein
series have coefficients $\sum_{d\mid n}\chi(d)$ and constant $L(0,\chi)/2$
\cite[\S4.8]{mDiamond}. By linearity,
$L(0,c_a)=-\sum_{d=1}^{5}dc_a(d)/5=\eta_a$, which gives exactly $f_a$.
Let $F=f_1-4f_2$. At a zero of $H$, one has $F=1/10$.
The modular function $J=f_2/F$ is nonconstant because $F$ has coefficient one
at $t$, $f_2$ has coefficient zero there and has nonzero constant term.

The modular-function field with algebraic Fourier coefficients is algebraic
over $\overline{\mathbb Q}(J)$, since $J$ is nonconstant on the modular
curve \cite{mDiamond,mBhowmik}. In particular the functions
\[
 \frac{E_4}{F^4},\qquad \frac{E_6}{F^6},\qquad
 \frac{\theta F-E_2F/12}{F^3}
\]
are algebraic over that field of rational functions in $J$. The third
numerator is the weight-one Serre derivative, which has weight three.
All functions displayed here are finite at the zero. Specializing a
primitive polynomial relation in $\overline{\mathbb Q}[J][X]$ leaves a
nonzero polynomial there, since its coefficients have no common root.
As $F=1/10$, the first two ratios put $E_4,E_6$ in an algebraic extension
of $\overline{\mathbb Q}(J)$ at the zero. The third puts $E_2$ in an
algebraic extension of $\overline{\mathbb Q}(J,\theta F)$.
Finally $J=10(Y_2-1/10)$, $\theta F=-5QH'(Q)$, and $t^5=1/Q$.
Nesterenko's bound therefore gives transcendence degree at least three for
the field generated by $Q,Y_2,H'(Q)$ at the zero.
\end{proof}

\subsection{Hyperbolic geometry and a zeta derivative}

\par\hypertarget{cat-M22}{}\label{page:M22}\hypertarget{cat-M23}{}\label{page:M23}
\begin{proposition}[M22--M23]
Let $\ell_n$ be the length of the branch geodesic in the $n$-fold cyclic
branched cover of the figure-eight knot with its hyperbolic metric, for
$n=5,7,11$. The ratios $\ell_5/\ell_7$ (M22) and $\ell_5/\ell_{11}$
(M23) are transcendental.
\end{proposition}
\begin{proof}
The branch geodesic maps isometrically to the singular locus of the cone
orbifold of angle $2\pi/n$. The figure-eight length formula
\cite{mMednykh} gives
\[
 \ell_n=2\log U_n,\qquad
 U_n=A_n+\sqrt{A_n^2-1},\qquad
 A_n=1+\cos(2\pi/n)-\cos(4\pi/n)>1 .
\]
For $p=5,7,11$ write $F_p=\mathbb Q(\cos(2\pi/p))$. Every
$U_p^m+U_p^{-m}=2T_m(A_p)$ for $m\ge1$ lies in $F_p$, where $T_m$ is
the Chebyshev polynomial. The identity embedding makes this number larger
than two. The automorphism carrying $\cos(2\pi/p)$ to
$c=\cos(\pi-\pi/p)$ sends $A_p$ to $2+c-2c^2$. Since $-1<c<-1/2$,
this lies in $(-1,1)$, so the conjugate trace has absolute value at most two.
The trace is therefore irrational. For $p=7$ or $11$ the degrees of $F_5$
and $F_p$ are coprime, so their intersection is $\mathbb Q$. A rational
length ratio would give $U_5^m=U_p^n$ for positive integers $m,n$ and hence
a common irrational trace, which is impossible. An algebraic irrational
ratio would make $U_5=U_p^{\ell_5/\ell_p}$ transcendental by the
Gelfond--Schneider theorem, again impossible.
\end{proof}

\par\hypertarget{cat-M24}{}\label{page:M24}\hypertarget{cat-M25}{}\label{page:M25}
\begin{proposition}[M24--M25]
The maximal cusp packing density $d$ of the Gieseking manifold (M24) and
the zeta derivative $\zeta'(-2)$ (M25) are irrational.
\end{proposition}
\begin{proof}
Write $\chi_{-3}$ for the nontrivial character modulo three. The regular
ideal tetrahedron giving the Gieseking manifold has volume
$v_3=3\sqrt3\,L(2,\chi_{-3})/4$ \cite{mCDT}. Normalize its vertices
in the upper half-space model to $0,1,e^{i\pi/3},\infty$. The four equally
sized cusp sectors attain tangency at height one at the vertex $\infty$,
with Euclidean cross-section an equilateral triangle of side one.
Its hyperbolic volume is
$(\sqrt3/4)\int_1^\infty z^{-3}\,dz=\sqrt3/8$. The single cusp identifies
the four sectors, so first tangency is the maximal cusp and its volume is
$\sqrt3/2$. Thus $d=\sqrt3/(2v_3)=2/(3L(2,\chi_{-3}))$.
Calegari, Dimitrov and Tang proved that $1,\zeta(2),L(2,\chi_{-3})$
are linearly independent over $\mathbb Q$ \cite{mCDT}. Rationality of
$d$ would make $L(2,\chi_{-3})$ rational.

For M25, differentiation of the Riemann zeta functional equation at the
simple zero of $\sin(\pi s/2)$ at $s=-2$ gives
\[
 \zeta'(-2)=-\frac{\zeta(3)}{4\pi^2}
           =-\frac{\zeta(3)}{24\zeta(2)}.
\]
Qian Tang's publicly released Lean proof \cite{mTang} establishes
$\mathbb Q$-linear independence of $1,\zeta(2),\zeta(3)$. A rational
value of the displayed quotient would give a nontrivial rational relation
between the last two numbers. The argument uses the theorem stated in that public formalization as an external input. Its Lean build was not replayed for this manuscript.
\end{proof}



\section{Arithmetic series, digital constants, and effective avoidance}
\label{sec:arith}
Some series in this section can be excluded from an explicit number field.
For others, a family of $E$-function values, a small integer linear form,
or the digits of an expansion give the required arithmetic conclusion.

\subsection{Six individual \texorpdfstring{$q$}{q}-Stieltjes values}
For $q>1$ and $x>0$, use the constant Laurent coefficient
$\gamma_0(q,x)$ at $s=1$ of
\[
 \zeta_q(s,x)=\sum_{m\ge0}
       \frac{q^{m+x}}{([m+x]_q)^s},\qquad
       [u]_q=\frac{q^u-1}{q-1}
       \quad(\Re s>1).
\]
Chatterjee and Garg's coefficient formula \cite[Eq.~(7)]{nCG} is
\[
 \gamma_0(q,x)=(q-1)
 \left(L_q(q^{1-x})+\frac{\log(q-1)}{\log q}+\frac12-x\right),
 \qquad L_q(z)=\sum_{m\ge1}\frac{z^m}{q^m-1}.
\]
The Lambert series converges for $|z|<q$.

\par\hypertarget{cat-N01}{}\label{page:N01}\hypertarget{cat-N02}{}\label{page:N02}\hypertarget{cat-N03}{}\label{page:N03}\hypertarget{cat-N04}{}\label{page:N04}\hypertarget{cat-N05}{}\label{page:N05}\hypertarget{cat-N06}{}\label{page:N06}
\begin{proposition}[N01--N06. Plastic-number specializations]
Let $\rho>1$ be the real root of $X^3-X-1$. Each entry below lies outside
$\mathbb Q(\rho)$ and hence is irrational.
\[
\begin{array}{c|rrrrrr}
\text{entry}&\mathrm{N01}&\mathrm{N02}&\mathrm{N03}&\mathrm{N04}&\mathrm{N05}&\mathrm{N06}\\
q&\rho&\rho^5&\rho^5&\rho^5&\rho^5&\rho^5\\
x&1&1/5&1&2/5&3/5&4/5
\end{array}
\]
\end{proposition}
\begin{proof}
For a Pisot number $q$ and a nonzero algebraic integer $\alpha\in K=\mathbb Q(q)$,
Dixit, Kumar and Pathak's Corollary~1.2 gives $L_q(\alpha)\notin K$
when the two inequalities below hold \cite{nDKP}.
\[
 |\alpha|<q,\qquad\min(1,|\alpha|)\prod_{\sigma:K\hookrightarrow\mathbb C}
\max(1,|\sigma(\alpha)|)<q.
\]
The cubic has no rational root and discriminant $-23$. Its conjugate pair
has common modulus $\rho^{-1/2}$ by the product of the three roots.
Thus $\rho$ is Pisot. Since $\rho^5=\rho^4+1$, both
$\rho^4(\rho-1)=1$ and $\rho=\rho^5/(\rho^5-1)$ hold. In particular
$\rho^5$ is Pisot with $\mathbb Q(\rho^5)=\mathbb Q(\rho)$.
At $q=\rho,x=1$ the Lambert argument is $1$, its height is $1<\rho$,
and the logarithmic ratio is $-4$. At $q=\rho^5,x=j/5$ it is
$\alpha=\rho^{5-j}$ with height $\rho^{5-j}<\rho^5$, for each
$j=1,\ldots,5$. The logarithmic ratio is $4/5$. Subtracting the
remaining rational term and dividing by $q-1$ would put the
corresponding Lambert value in $\mathbb Q(\rho)$ if the listed constant
belonged to that field.
\end{proof}

\subsection{Reciprocal recurrences and binary divisor sums}
Write $F_0=0,F_1=1$ and $L_0=2,L_1=1$. Both sequences satisfy
$u_{n+2}=u_{n+1}+u_n$. In this subsection $L_n$ is a Lucas number.

\par\hypertarget{cat-N07}{}\label{page:N07}\hypertarget{cat-N08}{}\label{page:N08}\hypertarget{cat-N09}{}\label{page:N09}\hypertarget{cat-N10}{}\label{page:N10}\hypertarget{cat-N11}{}\label{page:N11}
\begin{proposition}[N07--N11. Five reciprocal sums]
The three sums, in order,
\[
 \mathrm{N07}=\sum_{n\ge1}\frac1{F_n},\qquad
 \mathrm{N08}=\sum_{n\ge1}\frac1{F_{2n}},\qquad
 \mathrm{N09}=\sum_{n\ge1}\frac1{L_n}
\]
lie outside $\mathbb Q(\sqrt5)$. With $d(n)$ denoting the number of
positive divisors of $n$, the numbers
\[
 \mathrm{N10}=\sum_{n\ge1}\frac{d(n)}{2^n},\qquad
 \mathrm{N11}=\sum_{n\ge1}\frac{1}{2^n+1}
\]
are linearly independent together with $1$ over $\mathbb Q$.
\end{proposition}
\begin{proof}
Put $\varphi=(1+\sqrt5)/2$ and $\psi=-1/\varphi$.
For $R_n=g_1\varphi^n+g_2\psi^n$ with the Fibonacci or Lucas initial
conditions, Tachiya's recurrence criterion \cite[Cor.~2 and Ex.~2]{nTach}
puts $\sum_{n\ge1}1/R_{an+b}$ outside $\mathbb Q(\sqrt5)$ for
$a\ge1,b\ge0$. In these instances the coefficient ratio $g_1/g_2$ is
$-1$ or $1$, the characteristic-root ratio is a unit, and $\varphi^a$
is integral, so its denominator to the fourth power is $1<|\varphi/\psi|^a$.
The argument $1$ is inside the disk of convergence. Take
$(R,a,b)=(F,1,0),(F,2,0),(L,1,0)$.

Tachiya's periodic-coefficient theorem \cite[Thm.~1]{nTach} puts
$\sum a_n/(2^n-1)$ outside $\mathbb Q$ for every nonzero rational
sequence of period dividing two. Expanding the geometric series gives
\[
 \sum_{n\ge1}\frac1{2^n-1}=\sum_{m\ge1}d(m)2^{-m},\qquad
 \sum_{n\ge1}\frac{(-1)^{n-1}}{2^n-1}=\sum_{k\ge1}\frac1{2^k+1}.
\]
The rearrangements are absolutely convergent. A rational linear
combination with coefficients $u,v$ has weights $u+v(-1)^{n-1}$.
Their simultaneous vanishing requires $u=v=0$.
\end{proof}

\subsection{A small positive integer for Ap\'ery's constant}
\par\hypertarget{cat-N12}{}\label{page:N12}
\begin{proposition}[N12. Ap\'ery's constant]
The number $\zeta(3)=\sum_{m\ge1}m^{-3}$ is irrational.
\end{proposition}
\begin{proof}
Let $D=1-(1-xy)z$, $H=x(1-x)y(1-y)z(1-z)/D$, and integrate
$I_n=\int_{[0,1]^3}H^n/D\,dx\,dy\,dz$. Following Beukers
\cite{nBeuk}, put $P_n(x)=(d/dx)^n[x^n(1-x)^n]/n!$.
Its coefficients are integers, since the $n$th derivative of a monomial
divided by $n!$ has integral coefficients. Start with the triple integral
of $P_n(x)P_n(y)/D$. Integrating by parts $n$ times in $x$ replaces
$P_n(x)/D$ by $x^n(1-x)^n(yz)^n/D^{n+1}$. For fixed $x,y>0$ make
the substitution $z=(1-w)/(1-(1-xy)w)$. With
$D_w=1-(1-xy)w$, its Jacobian in absolute value is $xy/D_w^2$
and $D=xy/D_w$. The resulting integrand is
$(1-x)^n(1-w)^n P_n(y)/D_w$. Integrating by parts $n$ times in
$y$ recovers $I_n$ with $w$ as its third variable. Boundary terms vanish
because $x^n(1-x)^n$ and its first $n-1$ derivatives vanish at both endpoints.

Integrate the original kernel in $z$ to obtain $-\log(xy)/(1-xy)$.
For a monomial $x^r y^s$ in $P_n(x)P_n(y)$, its integral against this
kernel is $2\zeta(3)-2\sum_{j=1}^r j^{-3}$ if $r=s$.
For $r<s$ it is $(s-r)^{-1}\sum_{j=r+1}^s j^{-2}$, and the case
$s<r$ follows by symmetry. Indeed, expanding the geometric series in
$xy$ gives terms $(A^{-2}-B^{-2})/(s-r)$ with $A=k+r+1$,
$B=k+s+1$, which telescope. The case $r=s$ follows from
$\int_0^1 x^{A-1}(-\log x)\,dx=A^{-2}$. Since $r,s\le n$,
this proves $d_n^3 I_n\in\mathbb Z+\mathbb Z\zeta(3)$ for
$d_n=\operatorname{lcm}(1,\ldots,n)$. It also gives $I_0=2\zeta(3)$.

For $u=\sqrt{xy}$, $(1-x)(1-y)\le(1-u)^2$ and differentiating in $z$
gives $\max_{0\le z\le1}z(1-z)/D=1/(1+u)^2$ for $u>0$.
Maximizing $u(1-u)/(1+u)$ at $u=\sqrt2-1$ yields
$0<I_n\le 2\zeta(3)(\sqrt2-1)^{4n}$. By the prime number theorem
$d_n=\exp(n+o(n))$. If $v\zeta(3)$ were integral for a positive
integer $v$, the positive integers $vd_n^3I_n$ would tend to zero,
since $e^3(\sqrt2-1)^4<1$.
\end{proof}

\subsection{Euler--Gompertz combinations and logarithmic moments}
Let $\gamma=\lim_{n\to\infty}(\sum_{k=1}^n k^{-1}-\log n)$ and
$\delta=e\int_1^\infty e^{-t}\,dt/t$. For $j\ge1$ put
\[
 M_j=\frac{1}{(j-1)!}\int_0^1
         \frac{1-e^{-t}}t(-\log t)^{j-1}\,dt,\qquad
 C_r=\frac{e\gamma+\delta}{e+r}.
\]

\par\hypertarget{cat-N13}{}\label{page:N13}\hypertarget{cat-N14}{}\label{page:N14}\hypertarget{cat-N15}{}\label{page:N15}\hypertarget{cat-N16}{}\label{page:N16}\hypertarget{cat-N17}{}\label{page:N17}
\begin{proposition}[N13--N17. Five transcendental numbers]
$\mathrm{N13}=C_0$, $\mathrm{N14}=C_1$, $\mathrm{N15}=C_{-1}$,
$\mathrm{N16}=M_2$, and $\mathrm{N17}=M_3$ are transcendental.
Every pair among $C_0,C_1,C_{-1}$ is algebraically independent.
\end{proposition}
\begin{proof}
The identity $\Gamma'(1)=-\gamma$ and integration by parts give
\[
 -\gamma=\int_0^\infty e^{-t}\log t\,dt
 =-\int_0^1\frac{1-e^{-t}}t\,dt+
  \int_1^\infty\frac{e^{-t}}t\,dt .
\]
Thus $M_1=\gamma+\delta/e$. Termwise integration, justified by the
absolutely summable integrated terms, gives
$M_j=\sum_{k\ge1}(-1)^{k-1}/(k^j k!)$.
Shidlovskii's theorem for the functions
$w_j(z)=1+\sum_{k\ge1}z^k/(k^j k!)$ says that
$w_0(a),\ldots,w_m(a)$ are algebraically independent at every
nonzero algebraic $a$ \cite[Ch.~7, \S1, Thm.~1]{nShid}.
At $a=-1$ it gives algebraic independence of
$e,M_1,M_2,M_3$, by inversion of $e^{-1}$ and the affine
changes $1-M_j$. For distinct $r,s\in\{0,1,-1\}$, positivity of
$M_1$ gives $C_r\ne C_s$, and the identities
\[
 C_r=\frac{eM_1}{e+r},\qquad
 e=\frac{sC_s-rC_r}{C_r-C_s},\qquad M_1=C_r(1+r/e)
\]
show that $\mathbb Q(C_r,C_s)=\mathbb Q(e,M_1)$ has transcendence
degree two. The proof does not assign an arithmetic status to $\gamma$
or to $\delta$ individually.
\end{proof}

\subsection{A coprime square-shell expansion}
\par\hypertarget{cat-N18}{}\label{page:N18}
\begin{proposition}[N18. Binary coprime square-shell constant]
For every integer $b\ge2$ define
\[
 D_b=\sum_{n\ge1}b^{-n}\,\mathbf 1_{\gcd(n,\lfloor\sqrt n\rfloor)=1}.
\]
The archived $C_b$ is denoted $D_b$ here to distinguish it from $C_r$
above. The binary constant $\mathrm{N18}=D_2$ is irrational, as is
every $D_b$.
\end{proposition}
\begin{proof}
The digit at $n=k^2+r$, $0\le r\le 2k$, equals
$\mathbf 1_{\gcd(r,k)=1}$. For every $k\ge2$ there is a zero at $k^2$
and a one at $k^2+1$. These digits therefore give an unambiguous
infinite base-$b$ expansion. If it had eventual period $h\ge1$, choose
a prime $p>h$ with $p^2$ in the periodic part. The digit at $p^2$
is zero, while the one at $p^2+h$ equals one because $\gcd(h,p)=1$.
Every rational number has eventually periodic base-$b$ digits, by long
division with finitely many possible remainders.
\end{proof}

\subsection{An explicit continued fraction over the field of Euler's constant}
The digit rules below specify the particular number in the November~11
archive \cite{nKeio}. They matter because changing a tie rule can change
the real number. A closed interval's margin is its distance to $0$.
All interval arithmetic takes endpoint extrema, and powers use repeated
interval multiplication.

\paragraph{Coefficient enclosures.}
At level $s\ge1$ write $N=s+2$. Enclose $e$ by the partial factorial
sum through $N$ with error interval $[0,1/(NN!)]$.
Enclose $\pi$ by Machin's formula $16\arctan(1/5)-4\arctan(1/239)$,
using the first $N$ alternating terms at $x$ and symmetric error
$x^{2N+1}/(2N+1)$. For rational $1\le t\le2$ put $z=(t-1)/(t+1)$
and enclose $\log t$ by $2\sum_{j=0}^{N-1}z^{2j+1}/(2j+1)$
with error interval $[0,2z^{2N+1}/((2N+1)(1-z^2))]$.
If $N=2^m t$, $1\le t<2$, combine the intervals for
$m\log2+\log t$ to obtain $J_{\log N}$. For Bernoulli numbers
$B_k$ in the convention $B_1=-1/2$ and $H_N=\sum_{j=1}^N1/j$, use
\[
 J_\gamma=H_N-\frac1{2N}+\sum_{k=1}^{N-1}\frac{B_{2k}}{2kN^{2k}}
             -J_{\log N}+[-R,R],
 \qquad R=\frac{|B_{2N}|}{2N\,N^{2N}} .
\]
The positive-real digamma remainder is bounded by the first omitted
term \cite[\S5.11]{aDLMF}. This also follows by expanding
$1/(N^2+t^2)$ in Binet's integral with finite geometric remainder.
The identity $|B_{2N}|=2(2N)!\zeta(2N)/(2\pi)^{2N}$ shows $R\to0$.
All three coefficient intervals contain their targets and shrink to points.

\paragraph{Ordering and selection.}
For a term $cX^xE^yT^zG^g$ write $\alpha=(x,y,z,g)$.
List the nonzero terms by increasing $(|\alpha|,\alpha)$, with
$|\alpha|=x+y+z+g$. Enumerate the integer polynomials of positive
formal $X$-degree whose first coefficient is positive, first by
$\sum_{\alpha}(|c_\alpha|+1+|\alpha|)$ over nonzero terms, then
lexicographically by the lists of $(\alpha,c_\alpha)$. For $P_i$ so
obtained, let $P_i^*(X)=P_i(X,e,\pi,\gamma)$. Enclose each coefficient
by termwise substitution at level $s$, and denote by $V(P_i,s,I)$ the
interval obtained by Horner evaluation of the resulting polynomial
on $I$. Begin with the empty digit word $w$. Its cylinder is
$I(w)=[0,1]$. For a nonempty $w$ the cylinder endpoints are
$p/q$ and $(p+p')/(q+q')$, where $p/q$ and $p'/q'$ are the last two
convergents, with initial pairs $0/1,1/0$.

At stage $s$ consider all unprocessed $i\le s$. Activate an index
when some coefficient interval first has positive margin, and store
the margin $\eta_i$ of its first such coefficient in order of increasing
$X$-degree. If none is active, append $1$ and end this stage. Otherwise
select the least active unprocessed $i$. If $V(P_i,s,I(w))$ has
positive margin $A/B$ in lowest terms, append $1+((A+B+s)\bmod 2)$.
If its margin is zero, make the following extension.

\paragraph{Extension when a cylinder contains a possible zero.}
Write $d=\deg_X P_i$. For $0\le j\le d$ put $u_j=w1^j2$ and
$x_j=[0;u_j,\overline1]$. To enclose $x_j$ at precision $h$ let
$b_h=\lfloor\sqrt{5\cdot2^{2h}}\rfloor$ and pass the interval with
endpoints $(1+2^{-h}b_h)/2$, $(1+2^{-h}(b_h+1))/2$ through
$t\mapsto(pt+p')/(qt+q')$, using the convergents for $u_j$.
Denote the sorted image by $J_j(h)$. Choose the least $h\ge1$ at which
these intervals have pairwise distance exceeding $2^{-h}$. Set
$\lambda=\eta_i\,2^{-hd}/((d+1)2^d)$. Take the least $t\ge1$ for
which all $V(P_i,t,J_j(t))$ have width at most $\lambda/4$, then
choose the smallest $j$ for which the absolute midpoint exceeds
$\lambda/2$. Denote the degree-$k$ coefficient interval at level $t$
by $[a_k,b_k]$, and put $B=1+\sum_{k=1}^d k\max(|a_k|,|b_k|)$.
Take the least $r\ge0$ satisfying $4^r(\lambda/4)q_{u_j}^{\,2}>2B$,
where $q_{u_j}$ is the last convergent denominator for $u_j$.
Replace $w$ by $u_j1^{2r}$. Both branches mark $i$ as processed.
The eventual infinite word defines $\Theta=[0;w_\infty]$.

\par\hypertarget{cat-N19}{}\label{page:N19}
\begin{proposition}[N19. The Keio continued-fraction constant]
$\mathrm{N19}=\Theta$ is computable and transcendental over
$\mathbb Q(e,\pi,\gamma)$. Its irrationality exponent is two.
\end{proposition}
\begin{proof}
The $x_j$ are distinct because their infinite words differ. Thus the
separation search terminates. In the Lagrange interpolation basis for
these nodes, each numerator coefficient has modulus at most $2^d$ and
each denominator has modulus at least $2^{-hd}$. Since one coefficient
of $P_i^*$ has modulus at least $\eta_i$, some $|P_i^*(x_j)|\ge\lambda$.
The enclosures shrink at these points, so the midpoint test terminates
and the selected node has $|P_i^*(x_j)|>3\lambda/8$. The derivative
has absolute value below $B$ on $[0,1]$. Appending two ones at least
doubles a convergent denominator. The returned cylinder thus has
diameter at most $(4^rq_{u_j}^2)^{-1}<\lambda/(8B)$ and contains $x_j$.
Throughout it $|P_i^*|>\lambda/4$. The positive-margin branch also
certifies nonvanishing on its whole cylinder.

Every polynomial with nonzero specialization eventually has a
coefficient interval excluding zero. Once active it can be delayed
by at most $i-1$ smaller indices. The nested cylinders shrink
effectively to one number, since their denominators grow at least
as fast as Fibonacci numbers. Hence $P_i^*(\Theta)\ne0$ for all
nonzero specializations. Clearing denominators in an alleged
nonconstant relation over $\mathbb Q(e,\pi,\gamma)$ would produce
one of these polynomials, up to sign. Finally, the partial
quotients are at most two. The convergent formula gives
$|\Theta-p_n/q_n|>1/(4q_n^2)$, and Legendre's convergent criterion gives
$|\Theta-p/q|\ge1/(2q^2)$ for reduced nonconvergents. Infinitely
many convergents have error below $1/q_n^2$. The two bounds establish
the asserted exponent.
\end{proof}

\subsection{Three Tschakaloff values at the plastic number}
\par\hypertarget{cat-N20}{}\label{page:N20}\hypertarget{cat-N21}{}\label{page:N21}\hypertarget{cat-N22}{}\label{page:N22}
\begin{proposition}[N20--N22. A partial theta function and its derivatives]
For $q=\rho$, define $T_q(z)=\sum_{n\ge0}z^n q^{-n(n+1)/2}$.
The three numbers $\mathrm{N20}=T_{\rho}(1)$,
$\mathrm{N21}=T'_{\rho}(1)$, and $\mathrm{N22}=T''_{\rho}(1)$
are linearly independent together with $1$ over $\mathbb Q(\rho)$.
\end{proposition}
\begin{proof}
The series is entire by the ratio test, so the derivative entries
are the sums with weights $n$ and $n(n-1)$, respectively.
Dixit, Kumar and Pathak's Corollary~1.3 \cite{nDKP} gives independence of
$1,T_q(\alpha),\ldots,T_q^{(m)}(\alpha)$ over $\mathbb Q(q)$ under
the height condition used for N01. Take $q=\rho$, $\alpha=1$, $m=2$.
The conjugate product is $1<\rho$. If one of the three were in the
field it would yield a linear relation with $1$.
\end{proof}

\subsection{Prime-square indices and a quotient of logarithms}
\par\hypertarget{cat-N23}{}\label{page:N23}\hypertarget{cat-N24}{}\label{page:N24}
\begin{proposition}[N23--N24. Prime-square reciprocal sums]
The two numbers
\[
 \mathrm{N23}=\sum_{p\ {\rm prime}}\frac1{F_{p^2}},\qquad
 \mathrm{N24}=\sum_{p\ {\rm prime}}\frac1{L_{p^2}}
\]
are linearly independent together with $1$ over $\mathbb Q(\sqrt5)$.
\end{proposition}
\begin{proof}
The recurrence corollary of Duverney, Suzuki and Tachiya
\cite[Cor.~1.2]{nDST} applies to an increasing sequence of pairwise
coprime positive odd indices $m_i$ with $\sum_i1/m_i<\infty$.
For a real Pisot number $\alpha$, $\beta=\pm\alpha^{-1}$ and
$U_n=(\alpha^n-\beta^n)/(\alpha-\beta)$, $V_n=\alpha^n+\beta^n$,
it makes $1,\sum_i1/U_{m_i},\sum_i1/V_{m_i}$ linearly independent
over $\mathbb Q(\alpha)$. Take the squares of odd primes. Their
reciprocals sum to at most $\sum_{n\ge2}n^{-2}$. Taking
$\alpha=\varphi$ and $\beta=-\varphi^{-1}$ gives the Fibonacci
and Lucas sequences by their initial values. Adjoining the term $p=2$
adds $1/3$ and $1/7$ to the respective sums. Such rational
translations preserve independence together with $1$.
\end{proof}

\par\hypertarget{cat-N25}{}\label{page:N25}
\begin{proposition}[N25. The plastic and golden logarithms]
The number $\mathrm{N25}=\log\rho/\log\varphi$ is transcendental.
\end{proposition}
\begin{proof}
If the quotient were $m/n$ for positive integers $m,n$, the equality
$\rho^n=\varphi^m$ would give an element of
$\mathbb Q(\rho)\cap\mathbb Q(\sqrt5)=\mathbb Q$, because the degree
of the intersection divides both three and two. A nonreal conjugate
of $\rho$ has modulus $\rho^{-1/2}$, so its $n$th power cannot equal
$\rho^n>1$. This excludes rational exponents. If the quotient were
algebraic irrational, Gelfond--Schneider would make
$\varphi^{\,\log\rho/\log\varphi}$ transcendental, contradicting
its value $\rho$.
\end{proof}

\subsection{Rational Lambert bases in the research archive}
\label{sec:nlambertaudit}
The November~13 draft \cite{nLamDraft} proposes irrationality of
$F(q)=\sum_{n\ge1}(q^n-1)^{-1}$ at $q=21548/305$ and $20814/299$.
A separate checker \cite{nLamCode} reaches $16862/265$ as well.
Their numerical comparisons and the arithmetic implication can be
checked separately.

Write $q=a/b$ in lowest terms. The determinant argument for
dimensions $(h,k,L,s)$ requires positive evaluated
determinants $Q_N(q,F(q))$ and positive monic denominator clearers
$J_N(q)$, with $P_N=J_N Q_N\in\mathbb Z[U,X]$,
$\deg_U P_N\le d_N=\deg J_N$, and $\deg_X P_N\le hN$.
It further requires the estimates
\[
 \begin{aligned}
 N^{-3}\log Q_N(q,F(q))&\longrightarrow-S\log q,
 &\limsup_{N\to\infty}N^{-3}d_N&\le C,\\
 \log J_N(q)&=d_N\log q+o(N^3).
 \end{aligned}
\]
Under these hypotheses, if $F(q)=u/v$ with $v>0$, then
$b^{d_N}v^{hN}P_N(a/b,u/v)$ is a positive integer. Its logarithm
divided by $N^3$ has upper limit at most
$C\log a-S\log(a/b)$. Hence $\log b/\log a<1-C/S$ suffices to
contradict rationality.

For $(3000,161,3097,245)$ the supplied six-profile calculation yields
$A=2054711745$, $S=26766016500$. Re-evaluating the rational affine
pieces on $[3,3258]$ gives $4048$ positive pieces and, on summing
each integral rounded downward to a multiple of $10^{-9}$,
\[
 I\ \ge\ 6622155682389376084/10^9 .
\]
Here $I$ is the integral of $x$ times the proposed denominator saving.
With $C=A+L^3/8+(3h(k+L)^2-6I)/\pi^2$, the numerical expression
$1-C/S$ exceeds $0.57329026631440$. This is above both ratios in the
first draft. At $(100000,5378,103223,8156)$ the checker gives
$20369$ positive pieces on $[20,108601]$ and a lower cutoff bound
$0.5732902769252565$, above $\log265/\log16862$.
For reproduction, on each affine piece $[l,u]$ with saving
$\alpha+\beta x$ the integrated contribution is
$\alpha(u^2-l^2)/2+\beta(u^3-l^3)/3$.

The outstanding mathematical step is to prove that the six supplied
valuation expressions actually bound the poles of these specialized
determinants, including their changes under transposition and
prefix interchange, and hence give the required denominator clearer.
The cited development \cite{nLamSource} is the reference for that
transfer in the draft. The reproduced computations check the piecewise
integrals. The determinant identities for the new dimensions still require
verification. These three candidates therefore carry no numbered
irrationality conclusion in this paper.


\section{Recovery identities and unresolved arithmetic questions}\label{sec:open}
This section distinguishes an exact identity evaluated at a zero from the arithmetic conclusion one may wish to draw from it. None of the open targets below contributes to the count of one hundred.
\subsection{Euler's constant at a Bessel zero}\label{sec:euler}
Let $j$ be the first positive zero of $J_0$, and let $H_k=\sum_{r=1}^k1/r$. Put
\[
 B(z)=\sum_{k\ge1}\frac{(-1)^{k+1}H_k}{(k!)^2}(z^2/4)^k .
\]
The positive-real Neumann expansion for $Y_0$ reads
\[
 \frac{\pi}{2}Y_0(z)=(\gamma+\log(z/2))J_0(z)+B(z)
 \qquad(z>0)
\]
\cite[Eq.~10.8.2]{aDLMF}. The zero $j=2.4048255576\ldots$ is simple. It is transcendental by the Bessel theorem used for A15--A17. Differentiating before evaluating at $j$ yields the exact formula
\[
 \gamma=\frac{\frac\pi2Y_0'(j)-B'(j)}{J_0'(j)}-\log(j/2).
\]
In particular, direct substitution before differentiation would erase the $\gamma$ term. The displayed expression still involves the stated derivative values and the logarithm. No reduction of this expression to a nonconstant algebraic-coefficient rational function of $j$ is obtained here. N13--N17 prove transcendence for quantities involving $\gamma$ or its associated exponential integral, while the individual irrationalities of $\gamma$ and the Euler--Gompertz constant $\delta$ remain unresolved here.
\subsection{Odd zeta values recovered between Bessel zeros}\label{sec:odd}
For $z>0$ with $J_0(z)\ne0$, take logarithms of analytic germs at $\nu=0$, normalized to have constant term zero. For $m\ge2$ define
\begin{align*}
 L_m(z)&=[\nu^m]\log\frac{J_\nu(z)}{J_0(z)},\\
 B_m(z)&=[\nu^m]\log\!\left(\Gamma(1+\nu)(2/z)^\nu
                          \frac{J_\nu(z)}{J_0(z)}\right).
\end{align*}
Here $[\nu^m]$ extracts a Taylor coefficient, including the factorial in its denominator. For $m$ odd the log-gamma expansion gives $B_m=L_m-\zeta(m)/m$ \cite[\S5.7]{aDLMF}.
\begin{proposition}[A recovery point in each gap]\label{prop:odd}
For every odd integer $m\ge3$ and each pair of consecutive positive zeros $j_n,j_{n+1}$ of $J_0$, there exists $t\in(j_n,j_{n+1})$ such that $L_m(t)=0$. At every such point $\zeta(m)=-mB_m(t)$.
\end{proposition}
\begin{proof}
Put $a_k(z)=[\nu^k]J_\nu(z)/J_0(z)$ for $k\ge1$. The positive zeros of $J_0$ are simple, so each $a_k$ has at most a simple pole at any such $j$. The order derivative of $J_\nu$ at zero is $\pi Y_0/2$. Using $J_0Y_0'-J_0'Y_0=2/(\pi z)$ gives
\[
 a_1(z)=\frac{c_j}{z-j}+O(1),\qquad c_j=-\frac{1}{jJ_0'(j)^2}<0
\]
\cite[\S\S10.5, 10.15, 10.21]{aDLMF}. The only term of the logarithm contributing a pole of order $m$ is $(-1)^{m-1}a_1^m/m$. For odd $m$, $L_m$ consequently tends to $-\infty$ on the right of $j_n$ and $+\infty$ on the left of $j_{n+1}$. Between them it is real and continuous. The intermediate value theorem supplies a zero. The coefficient identity gives the recovery formula.
\end{proof}
For $m=3$ the first gap contains the numerically observed root near $3.805302720253$. The proposition requires neither uniqueness nor a decimal isolation of this root. It supplies no arithmetic classification of $B_m(t)$; the irrationality of $\zeta(3)$ in N12 uses a separate argument. In particular, its transcendence remains unresolved here.
\subsection{A criterion for the factorial-minus-one series}\label{sec:erdos}
The question whether $S=\sum_{n\ge2}(n!-1)^{-1}$ is irrational is Erd\H{o}s problem 68 \cite{pErdos}. Finite rational sums encode an equivalent question.
\begin{proposition}
Put $U_N=N!\sum_{n=2}^N1/(n!-1)$ and write $\{x\}=x-\lfloor x\rfloor$. Then $S$ is rational if and only if
\[
 1-\frac1N<\{U_N\}<1-\frac1{N+1}
 \qquad\text{for every sufficiently large }N .
\]
\end{proposition}
\begin{proof}
Set $V_N=N!S-U_N$. Its first term exceeds $1/(N+1)$. For $N\ge3$, replacing the successive factorial-tail ratios after the first by $1/(N+2)$ gives
\[
 V_N\le \frac{1}{1-1/(N+1)!}\frac{N+2}{(N+1)^2}<\frac1N,
\]
where the last step uses $(N+1)!>(N+1)^2$. If $S$ is rational then $N!S$ is eventually an integer, giving the stated condition. Conversely, eventual membership furnishes integers $a_N$ with $|N!S-a_N|<1/[N(N+1)]$. Consequently the integer $a_{N+1}-(N+1)a_N$ has absolute value at most $1/N+1/[(N+1)(N+2)]<1$ for all large $N$. Thus $a_N/N!$ is eventually a fixed rational, and its limit is $S$.
\end{proof}
An irrationality proof by this criterion requires infinitely many violations. Finite computations, of any prescribed length, leave that assertion open. The analogous analytic gap for the two rational-base Lambert values in the earlier chronology is given explicitly at the end of Section~\ref{sec:arith}.


\clearpage\appendix
\section{Concordance with the forty-three earlier entries}\label{sec:archive}
Rows follow the original research chronology. Times are Korea Standard Time. Entries sharing a minute retain their recorded order; a date alone means that the time was unspecified. The last column gives the proof or the precise remaining question.
\par\medskip
{\small\renewcommand{\arraystretch}{1.19}\setlength{\tabcolsep}{4pt}
\begin{longtable}{@{}r>{\raggedright\arraybackslash}p{27mm}>{\raggedright\arraybackslash}p{39mm}>{\raggedright\arraybackslash}p{77mm}@{}}
\toprule Row & 2026 (KST) & Archived subject & Result in this manuscript\\\midrule\endfirsthead
\toprule Row & 2026 (KST) & Archived subject & Result in this manuscript\\\midrule\endhead
\bottomrule\endfoot
01 & Oct 11 & Beta median & \hyperlink{cat-A01}{A01}. Median of $\operatorname{Beta}(3/2,5/2)$. \\
02 & Oct 11 & Gaussian mixture median & \hyperlink{cat-A05}{A05}. Unequal mixture on a bounded interval. \\
03 & Oct 11 & Truncated Gaussian median & \hyperlink{cat-A06}{A06}. Gaussian on $[0,1]$. \\
04 & Oct 11 & Exponential-integral root & \hyperlink{cat-A07}{A07}. The unique real solution with $t>1$. \\
05 & Oct 11 & Complementary error function & \hyperlink{cat-A08}{A08}. All zeros; upper incomplete gamma extension follows A08. \\
06 & Oct 17 & Modified Bessel extrema & \hyperlink{cat-A09}{A09--A10}. Both the maximizing argument and the maximum. \\
07 & Nov 11 & Coprime square shells & \hyperlink{cat-N18}{N18}. Irrationality for every integer base $b\ge2$. \\
08 & Nov 11 & Keio continued fraction & \hyperlink{cat-N19}{N19}. Transcendence over $\mathbb Q(e,\pi,\gamma)$; irrationality exponent two. \\
09 & Nov 13, 09:19 & Two rational Lambert bases & \hyperref[sec:nlambertaudit]{Lambert audit}. Two proposed cases. The denominator estimate remains to be justified. \\
10 & Nov 13, 21:47 & Laplace limit & \hyperlink{cat-A11}{A11}. Transcendence of the Kepler limit. \\
11 & Nov 13, 21:51 & Euler's constant at a Bessel zero & \hyperref[sec:euler]{Euler recovery}. Exact recovery at the first zero of $J_0$. The status of $\gamma$ remains open. \\
12 & Nov 13, 21:51 & Euler--Gompertz combinations & \hyperlink{cat-N13}{N13--N15}. Each pair is algebraically independent. \\
13 & Nov 13, 21:57 & Reciprocal Fibonacci sum & \hyperlink{cat-N07}{N07}. Outside $\mathbb Q(\sqrt5)$; transcendence unresolved here. \\
14 & Nov 13, 22:00 & Twice the Dawson maximum & \hyperlink{cat-A14}{A14}. The doubled maximum equals the reciprocal of its location. \\
15 & Nov 13, 22:00 & Bratu parameter & \hyperlink{cat-A12}{A12}. Unit interval; shares its root with A11. \\
16 & Nov 13, 22:03 & Square elliptic lattice & \hyperlink{cat-A18}{A18}. An algebraic nonzero multiple of a period. \\
17 & Nov 13, 22:07 & Bessel recovery of $\zeta(3)$ & \hyperref[sec:odd]{Odd zeta recovery}. An exact identity in every positive Bessel zero gap. \\
18 & Nov 13, 23:08 & Cubic P\'olya return probability & \hyperlink{cat-M04}{M04}. Independence with $\pi$ and $e^{\pi\sqrt6}$. \\
19 & Nov 13, 23:08 & Lattice resistance at $(1,1,0)$ & \hyperlink{cat-M05}{M05}. First of an algebraically independent pair. \\
20 & Nov 13, 23:08 & Lattice resistance at $(1,1,1)$ & \hyperlink{cat-M06}{M06}. Second resistance; joint relation with the return probability. \\
21 & Nov 13, 23:26 & Ap\'ery's constant & \hyperlink{cat-N12}{N12}. Positive linear forms prove irrationality. \\
22 & Nov 13, 23:27 & Upper Wythoff compositions & \hyperlink{cat-M09}{M09}. Parts $\lfloor k\varphi^2\rfloor$. \\
23 & Nov 13, 23:28:50 & Zeta derivative at $-2$ & \hyperlink{cat-M25}{M25}. Irrationality via a publicly released Lean proof. \\
24 & Nov 13, 23:32 & Euler--Gompertz constant & \hyperref[sec:euler]{Euler recovery}. Individual irrationality unresolved here. Related combinations in N13--N15. \\
25 & Nov 13, 23:36 & Four-dimensional catenoid & \hyperlink{cat-A19}{A19}. Critical separation divided by boundary radius. \\
26 & Nov 13, 23:36 & Annular Steklov eigenvalue & \hyperlink{cat-A13}{A13}. The square of the normalized eigenvalue is irrational. \\
27 & Nov 13, 23:47 & Random $3$-XORSAT & \hyperlink{cat-G11}{G11}. Threshold in equations per variable. \\
28 & Nov 13, 23:52 & Factorial-minus-one sum & \hyperref[sec:erdos]{Factorial criterion}. Eventual fractional-part criterion. Irrationality unresolved. \\
29 & Nov 13, 23:54 & Figure-eight branched covers & \hyperlink{cat-M22}{M22}. The ratio of branch lengths at degrees five and seven. \\
30 & Nov 13, 23:54 & Gieseking cusp packing & \hyperlink{cat-M24}{M24}. Irrationality of the density; the volume itself is not classified here. \\
31 & Nov 13, 23:54 & Airy half-energy radius & \hyperlink{cat-A15}{A15}. Dimensionless unobscured pupil. \\
32 & Nov 13, 23:59 & Obscured half-energy radius & \hyperlink{cat-A16}{A16}. Central obscuration ratio one half. \\
33 & Nov 14, 00:14 & Discriminant maximum point & \hyperlink{cat-M01}{M01}. Unique maximizing positive nome. \\
34 & Nov 14, 00:14 & Discriminant maximum value & \hyperlink{cat-M02}{M02}. Independent of its location and curvature. \\
35 & Nov 14, 00:19 & Mixed Mahler growth & \hyperlink{cat-M13}{M13}. Zero- and one-run lengths from powers of two and three. \\
36 & Nov 14, 00:20 & Generalized pentagonal growth & \hyperlink{cat-M16}{M16}. Growth rate and leading coefficient are algebraically independent. \\
37 & Nov 14, 00:20 & Teardrop elastica parameter & \hyperlink{cat-A20}{A20}. Parameter $m_T$ in the Legendre convention. \\
38 & Nov 14, 00:20 & Heart elastica parameter & \hyperlink{cat-A21}{A21}. Parameter $m_H$ in the same convention. \\
39 & Nov 14, 00:20 & Elastica length ratio & \hyperlink{cat-A22}{A22}. Teardrop length divided by the wave half-period. \\
40 & Nov 14, 00:23 & Reflected $q$-Stieltjes zero & \hyperlink{cat-M19}{M19--M21}. The root, a reflected difference and a derivative are algebraically independent. \\
41 & Nov 14, 00:58 & Plastic $q$-Stieltjes value & \hyperlink{cat-N01}{N01}. $\gamma_0(\rho,1)\notin\mathbb Q(\rho)$. \\
42 & Nov 14, 00:58 & Fifth-power base, shift $1/5$ & \hyperlink{cat-N02}{N02}. $\gamma_0(\rho^5,1/5)\notin\mathbb Q(\rho)$. \\
43 & Nov 14, 00:58 & Fifth-power base, shift $1$ & \hyperlink{cat-N03}{N03}. $\gamma_0(\rho^5,1)\notin\mathbb Q(\rho)$. \\
\end{longtable}}
\noindent The graph $3$-core threshold added after this chronology is G01. Its associated birth fraction is G05. The $3$-uniform $2$-core edge density is G08 and the $3$-XORSAT threshold is G11.

\clearpage\section{Index of the one hundred quantities}\label{sec:directory}
The identifiers link to the proposition or family theorem giving the proof. The adjacent page number locates its statement in a printed copy.
\subsection*{Graph and probability models (G) and analytic quantities (A)}
All entries below are proved transcendental except A13 and A22, for which irrationality is proved. A13 also has an irrational square.
\par\medskip
\noindent{\small\setlength{\tabcolsep}{3pt}\renewcommand{\arraystretch}{1.50}%
\begin{tabularx}{\textwidth}{@{}l>{\raggedright\arraybackslash}Xrl>{\raggedright\arraybackslash}Xr@{}}\toprule
ID & Quantity & Page & ID & Quantity & Page \\\midrule
\hyperlink{cat-G01}{\textbf{G01}} & Graph $3$-core threshold & \pageref*{page:G01} & \hyperlink{cat-A01}{\textbf{A01}} & Beta$(3/2,5/2)$ median & \pageref*{page:A01} \\
\hyperlink{cat-G02}{\textbf{G02}} & Graph $4$-core threshold & \pageref*{page:G02} & \hyperlink{cat-A02}{\textbf{A02}} & Beta$(3/2,5/2)$ 2.5th percentile & \pageref*{page:A02} \\
\hyperlink{cat-G03}{\textbf{G03}} & Graph $5$-core threshold & \pageref*{page:G03} & \hyperlink{cat-A03}{\textbf{A03}} & Beta$(3/2,5/2)$ 97.5th percentile & \pageref*{page:A03} \\
\hyperlink{cat-G04}{\textbf{G04}} & Graph $6$-core threshold & \pageref*{page:G04} & \hyperlink{cat-A04}{\textbf{A04}} & Student $t_3$ 97.5th percentile & \pageref*{page:A04} \\
\hyperlink{cat-G05}{\textbf{G05}} & Newborn $3$-core vertex fraction & \pageref*{page:G05} & \hyperlink{cat-A05}{\textbf{A05}} & Truncated Gaussian mixture median & \pageref*{page:A05} \\
\hyperlink{cat-G06}{\textbf{G06}} & Newborn $4$-core vertex fraction & \pageref*{page:G06} & \hyperlink{cat-A06}{\textbf{A06}} & Gaussian median on $[0,1]$ & \pageref*{page:A06} \\
\hyperlink{cat-G07}{\textbf{G07}} & Newborn $5$-core vertex fraction & \pageref*{page:G07} & \hyperlink{cat-A07}{\textbf{A07}} & Positive $\operatorname{Ein}$ root above one & \pageref*{page:A07} \\
\hyperlink{cat-G08}{\textbf{G08}} & $3$-uniform $2$-core edge density & \pageref*{page:G08} & \hyperlink{cat-A08}{\textbf{A08}} & First upper-half-plane erfc zero & \pageref*{page:A08} \\
\hyperlink{cat-G09}{\textbf{G09}} & $4$-uniform $2$-core edge density & \pageref*{page:G09} & \hyperlink{cat-A09}{\textbf{A09}} & Modified Bessel maximum location & \pageref*{page:A09} \\
\hyperlink{cat-G10}{\textbf{G10}} & $3$-uniform $3$-core edge density & \pageref*{page:G10} & \hyperlink{cat-A10}{\textbf{A10}} & Modified Bessel maximum value & \pageref*{page:A10} \\
\hyperlink{cat-G11}{\textbf{G11}} & $3$-XORSAT threshold & \pageref*{page:G11} & \hyperlink{cat-A11}{\textbf{A11}} & Laplace limit of Kepler's equation & \pageref*{page:A11} \\
\hyperlink{cat-G12}{\textbf{G12}} & $4$-XORSAT threshold & \pageref*{page:G12} & \hyperlink{cat-A12}{\textbf{A12}} & Unit-interval Bratu critical parameter & \pageref*{page:A12} \\
\hyperlink{cat-G13}{\textbf{G13}} & $5$-XORSAT threshold & \pageref*{page:G13} & \hyperlink{cat-A13}{\textbf{A13}} & Annular Steklov eigenvalue $S$ & \pageref*{page:A13} \\
\hyperlink{cat-G14}{\textbf{G14}} & Orientation, 3 choices and capacity 2 & \pageref*{page:G14} & \hyperlink{cat-A14}{\textbf{A14}} & Twice the Dawson maximum & \pageref*{page:A14} \\
\hyperlink{cat-G15}{\textbf{G15}} & Orientation, 4 choices and capacity 2 & \pageref*{page:G15} & \hyperlink{cat-A15}{\textbf{A15}} & Airy half-energy radius & \pageref*{page:A15} \\
\hyperlink{cat-G16}{\textbf{G16}} & Orientation, 3 choices and capacity 3 & \pageref*{page:G16} & \hyperlink{cat-A16}{\textbf{A16}} & Annular half-energy radius ($h=1/2$) & \pageref*{page:A16} \\
\hyperlink{cat-G17}{\textbf{G17}} & Binary search tree height coefficient & \pageref*{page:G17} & \hyperlink{cat-A17}{\textbf{A17}} & Airy first dark-ring radius & \pageref*{page:A17} \\
\hyperlink{cat-G18}{\textbf{G18}} & Binary search tree saturation coefficient & \pageref*{page:G18} & \hyperlink{cat-A18}{\textbf{A18}} & Square elliptic-lattice sum & \pageref*{page:A18} \\
\hyperlink{cat-G19}{\textbf{G19}} & Poisson$(2)$ survival probability & \pageref*{page:G19} & \hyperlink{cat-A19}{\textbf{A19}} & Four-dimensional catenoid gap ratio & \pageref*{page:A19} \\
\hyperlink{cat-G20}{\textbf{G20}} & Gamma$(2,1)$ median & \pageref*{page:G20} & \hyperlink{cat-A20}{\textbf{A20}} & Teardrop elastica parameter & \pageref*{page:A20} \\
\hyperlink{cat-G21}{\textbf{G21}} & Gamma$(3,1)$ median & \pageref*{page:G21} & \hyperlink{cat-A21}{\textbf{A21}} & Heart elastica parameter & \pageref*{page:A21} \\
\hyperlink{cat-G22}{\textbf{G22}} & Gamma$(5,1)$ 95th percentile & \pageref*{page:G22} & \hyperlink{cat-A22}{\textbf{A22}} & Teardrop arc-length ratio & \pageref*{page:A22} \\
\hyperlink{cat-G23}{\textbf{G23}} & Planck peak per unit frequency & \pageref*{page:G23} & \hyperlink{cat-A23}{\textbf{A23}} & Clamped--free first beam root & \pageref*{page:A23} \\
\hyperlink{cat-G24}{\textbf{G24}} & Planck peak per unit wavelength & \pageref*{page:G24} & \hyperlink{cat-A24}{\textbf{A24}} & Clamped--clamped first beam root & \pageref*{page:A24} \\
\hyperlink{cat-G25}{\textbf{G25}} & Planck peak per log frequency & \pageref*{page:G25} & \hyperlink{cat-A25}{\textbf{A25}} & First positive sinc sidelobe location & \pageref*{page:A25} \\
\bottomrule\end{tabularx}}
\clearpage
\subsection*{Modular and geometric quantities (M) and arithmetic series (N)}
M01--M23, N13--N17, N19 and N25 are proved transcendental. The remaining entries below are proved irrational. Several propositions give the stronger conclusion of linear independence over a specified field.
\par\medskip
\noindent{\small\setlength{\tabcolsep}{3pt}\renewcommand{\arraystretch}{1.50}%
\begin{tabularx}{\textwidth}{@{}l>{\raggedright\arraybackslash}Xrl>{\raggedright\arraybackslash}Xr@{}}\toprule
ID & Quantity & Page & ID & Quantity & Page \\\midrule
\hyperlink{cat-M01}{\textbf{M01}} & Nome maximizing the discriminant & \pageref*{page:M01} & \hyperlink{cat-N01}{\textbf{N01}} & Plastic base $\gamma_0(\rho,1)$ & \pageref*{page:N01} \\
\hyperlink{cat-M02}{\textbf{M02}} & Maximum of the discriminant & \pageref*{page:M02} & \hyperlink{cat-N02}{\textbf{N02}} & Fifth power $\gamma_0(\rho^5,1/5)$ & \pageref*{page:N02} \\
\hyperlink{cat-M03}{\textbf{M03}} & Negative second derivative at the peak & \pageref*{page:M03} & \hyperlink{cat-N03}{\textbf{N03}} & Fifth power $\gamma_0(\rho^5,1)$ & \pageref*{page:N03} \\
\hyperlink{cat-M04}{\textbf{M04}} & Simple cubic walk return probability & \pageref*{page:M04} & \hyperlink{cat-N04}{\textbf{N04}} & Fifth power $\gamma_0(\rho^5,2/5)$ & \pageref*{page:N04} \\
\hyperlink{cat-M05}{\textbf{M05}} & Lattice resistance to $(1,1,0)$ & \pageref*{page:M05} & \hyperlink{cat-N05}{\textbf{N05}} & Fifth power $\gamma_0(\rho^5,3/5)$ & \pageref*{page:N05} \\
\hyperlink{cat-M06}{\textbf{M06}} & Lattice resistance to $(1,1,1)$ & \pageref*{page:M06} & \hyperlink{cat-N06}{\textbf{N06}} & Fifth power $\gamma_0(\rho^5,4/5)$ & \pageref*{page:N06} \\
\hyperlink{cat-M07}{\textbf{M07}} & Face-centred cubic return probability & \pageref*{page:M07} & \hyperlink{cat-N07}{\textbf{N07}} & Reciprocal Fibonacci sum & \pageref*{page:N07} \\
\hyperlink{cat-M08}{\textbf{M08}} & Body-centred cubic return probability & \pageref*{page:M08} & \hyperlink{cat-N08}{\textbf{N08}} & Even-index reciprocal Fibonacci sum & \pageref*{page:N08} \\
\hyperlink{cat-M09}{\textbf{M09}} & Upper Wythoff composition growth & \pageref*{page:M09} & \hyperlink{cat-N09}{\textbf{N09}} & Reciprocal Lucas sum & \pageref*{page:N09} \\
\hyperlink{cat-M10}{\textbf{M10}} & Lower Wythoff composition growth & \pageref*{page:M10} & \hyperlink{cat-N10}{\textbf{N10}} & Binary divisor Lambert sum & \pageref*{page:N10} \\
\hyperlink{cat-M11}{\textbf{M11}} & Beatty growth, slope $1+\sqrt2$ & \pageref*{page:M11} & \hyperlink{cat-N11}{\textbf{N11}} & Sum of $1/(2^n+1)$ & \pageref*{page:N11} \\
\hyperlink{cat-M12}{\textbf{M12}} & Beatty growth, slope $\pi$ & \pageref*{page:M12} & \hyperlink{cat-N12}{\textbf{N12}} & Ap\'ery's constant $\zeta(3)$ & \pageref*{page:N12} \\
\hyperlink{cat-M13}{\textbf{M13}} & Mixed run-length growth, bases 2, 3 & \pageref*{page:M13} & \hyperlink{cat-N13}{\textbf{N13}} & Euler--Gompertz combination $C_0$ & \pageref*{page:N13} \\
\hyperlink{cat-M14}{\textbf{M14}} & Mixed run-length growth, bases 2, 5 & \pageref*{page:M14} & \hyperlink{cat-N14}{\textbf{N14}} & Euler--Gompertz combination $C_1$ & \pageref*{page:N14} \\
\hyperlink{cat-M15}{\textbf{M15}} & Mixed run-length growth, bases 2, 7 & \pageref*{page:M15} & \hyperlink{cat-N15}{\textbf{N15}} & Euler--Gompertz combination $C_{-1}$ & \pageref*{page:N15} \\
\hyperlink{cat-M16}{\textbf{M16}} & Generalized pentagonal growth & \pageref*{page:M16} & \hyperlink{cat-N16}{\textbf{N16}} & Second logarithmic moment $M_2$ & \pageref*{page:N16} \\
\hyperlink{cat-M17}{\textbf{M17}} & Square-part composition growth & \pageref*{page:M17} & \hyperlink{cat-N17}{\textbf{N17}} & Third logarithmic moment $M_3$ & \pageref*{page:N17} \\
\hyperlink{cat-M18}{\textbf{M18}} & Triangular-part composition growth & \pageref*{page:M18} & \hyperlink{cat-N18}{\textbf{N18}} & Binary coprime square-shell sum & \pageref*{page:N18} \\
\hyperlink{cat-M19}{\textbf{M19}} & Reflected $q$-Stieltjes zero $Q_*$ & \pageref*{page:M19} & \hyperlink{cat-N19}{\textbf{N19}} & Keio continued fraction $\Theta$ & \pageref*{page:N19} \\
\hyperlink{cat-M20}{\textbf{M20}} & Reflected difference $Y_2(Q_*)$ & \pageref*{page:M20} & \hyperlink{cat-N20}{\textbf{N20}} & Tschakaloff value $T_\rho(1)$ & \pageref*{page:N20} \\
\hyperlink{cat-M21}{\textbf{M21}} & Derivative $H'(Q_*)$ at the zero & \pageref*{page:M21} & \hyperlink{cat-N21}{\textbf{N21}} & First derivative $T'_\rho(1)$ & \pageref*{page:N21} \\
\hyperlink{cat-M22}{\textbf{M22}} & Figure-eight cover lengths $\ell_5/\ell_7$ & \pageref*{page:M22} & \hyperlink{cat-N22}{\textbf{N22}} & Second derivative $T''_\rho(1)$ & \pageref*{page:N22} \\
\hyperlink{cat-M23}{\textbf{M23}} & Figure-eight cover lengths $\ell_5/\ell_{11}$ & \pageref*{page:M23} & \hyperlink{cat-N23}{\textbf{N23}} & Prime-square reciprocal Fibonacci sum & \pageref*{page:N23} \\
\hyperlink{cat-M24}{\textbf{M24}} & Gieseking cusp packing density & \pageref*{page:M24} & \hyperlink{cat-N24}{\textbf{N24}} & Prime-square reciprocal Lucas sum & \pageref*{page:N24} \\
\hyperlink{cat-M25}{\textbf{M25}} & Zeta derivative $\zeta'(-2)$ & \pageref*{page:M25} & \hyperlink{cat-N25}{\textbf{N25}} & Plastic/golden logarithm ratio & \pageref*{page:N25} \\
\bottomrule\end{tabularx}}
\clearpage\phantomsection\label{sec:ref}\pdfbookmark[1]{References}{all-references}
\begin{thebibliography}{99}
\small
\bibitem{gPSW} B. Pittel, J. Spencer and N. Wormald,
Sudden emergence of a giant $k$-core in a random graph,
\emph{J. Combin. Theory Ser. B} \textbf{67} (1996), 111--151.
\href{https://cs.nyu.edu/~spencer/papers/k-core.pdf}{Author's paper}.

\bibitem{gMolloy} M. Molloy, Cores in random hypergraphs and Boolean formulas,
\emph{Random Structures Algorithms} \textbf{27} (2005), 124--135.
\href{https://doi.org/10.1002/rsa.20061}{doi:10.1002/rsa.20061}.

\bibitem{gCuckoo} M. Dietzfelbinger, A. Goerdt, M. Mitzenmacher, A. Montanari,
R. Pagh and M. Rink, Tight thresholds for cuckoo hashing via XORSAT,
\emph{ICALP 2010}, Part I, 213--225.
\href{https://arxiv.org/abs/0912.0287}{arXiv:0912.0287}.

\bibitem{gXor} B. Pittel and G. B. Sorkin, The satisfiability threshold for
$k$-XORSAT, \emph{Combin. Probab. Comput.} \textbf{25} (2016), 236--268.
\href{https://doi.org/10.1017/S0963548315000097}{doi:10.1017/S0963548315000097}.

\bibitem{gFKP} N. Fountoulakis, M. Khosla and K. Panagiotou,
The multiple-orientability thresholds for random hypergraphs,
\emph{Combin. Probab. Comput.} \textbf{25} (2016), 870--908.
\href{https://arxiv.org/abs/1309.6772}{arXiv:1309.6772}.

\bibitem{gDev} L. Devroye, A note on the height of binary search trees,
\emph{J. ACM} \textbf{33} (1986), 489--498.
\href{https://luc.devroye.org/devroye_1986_univ_a_note_on_the_height_of_binary_search_trees.pdf}{Author's paper}.

\bibitem{aOst} A. Ostrowski, Sur les relations alg\'ebriques entre les int\'egrales ind\'efinies, \emph{Acta Math.} \textbf{78} (1946), 315--318.

\bibitem{aBeu} F. Beukers, A refined version of the Siegel--Shidlovskii theorem, \emph{Ann. of Math.} \textbf{163} (2006), 369--379. \href{https://annals.math.princeton.edu/2006/163-1/p08}{Journal article}.

\bibitem{aDLMF} F. W. J. Olver et al., eds., \emph{NIST Digital Library of Mathematical Functions}. Chapters \href{https://dlmf.nist.gov/5}{5 (Gamma)}, \href{https://dlmf.nist.gov/7}{7 (Error functions)} and \href{https://dlmf.nist.gov/10}{10 (Bessel functions)}.

\bibitem{aSiegel} C. L. Siegel, \emph{Transcendental Numbers}, Annals of Mathematics Studies 16, Princeton University Press, 1949.

\bibitem{aFeld} N. I. Fel'dman and A. B. Shidlovskii, The development and present state of the theory of transcendental numbers, \emph{Russian Math. Surveys} \textbf{22} (1967), no.~3, 1--79. \href{https://doi.org/10.1070/RM1967v022n03ABEH001219}{Journal article}.

\bibitem{aMcM} C. T. McMullen, \emph{Differential and algebraic geometry of earthquakes}, preprint, 2026, \S12. \href{https://people.math.harvard.edu/~ctm/papers/home/text/papers/eq/eq.pdf}{Author's manuscript}.

\bibitem{aBuck} R. Buckmire, \emph{On exact and numerical solutions of the one-dimensional planar Bratu problem}, Occidental College, 2003. \href{https://sites.oxy.edu/ron/research/bratu/bratu.pdf}{Author's manuscript}.

\bibitem{aFraser} A. Fraser and R. Schoen, Sharp eigenvalue bounds and minimal surfaces in the ball, \emph{Invent. Math.} \textbf{203} (2016), 823--890. \href{https://arxiv.org/abs/1209.3789}{arXiv:1209.3789}.

\bibitem{aNest} Yu. V. Nesterenko, Modular functions and transcendence questions, \emph{Sb. Math.} \textbf{187} (1996), 1319--1348. \href{https://doi.org/10.4213/sm158}{Journal article}.

\bibitem{aPupil} J. J. Stamnes, H. Heier and S. Ljunggren, Encircled energy for systems with centrally obscured circular pupils, \emph{Appl. Opt.} \textbf{21} (1982), 1628--1633. \href{https://doi.org/10.1364/AO.21.001628}{Journal article}.

\bibitem{aSch} T. Schneider, Arithmetische Untersuchungen elliptischer Integrale, \emph{Math. Ann.} \textbf{113} (1937), 1--13.

\bibitem{aWald} M. Waldschmidt, Elliptic functions and transcendence, in K. Alladi, ed., \emph{Surveys in Number Theory}, Developments in Mathematics 17, Springer, 2008, 143--188. \href{https://webusers.imj-prg.fr/~michel.waldschmidt/articles/pdf/SurveyTrdceEllipt2006.pdf}{Author's manuscript}.

\bibitem{aMMR} T. Miura, M. M\"uller and F. Rupp, Optimal thresholds for preserving embeddedness of elastic flows, \emph{Amer. J. Math.} \textbf{147} (2025), 33--80. \href{https://doi.org/10.1353/ajm.2025.a950273}{Journal article}.

\bibitem{aSil} J. H. Silverman, \emph{The Arithmetic of Elliptic Curves}, 2nd ed., Graduate Texts in Mathematics 106, Springer, 2009.

\bibitem{mGlasser} M.~L. Glasser and J. Boersma, Exact values for the cubic lattice Green functions,
\emph{J. Phys. A} \textbf{33} (2000), 5017--5023.
\href{https://doi.org/10.1088/0305-4470/33/28/306}{DOI}.

\bibitem{mGutt} A.~J. Guttmann, Lattice Green's functions in all dimensions,
\emph{J. Phys. A} \textbf{43} (2010), 305205.
\href{https://doi.org/10.1088/1751-8113/43/30/305205}{DOI}.

\bibitem{mBeatty} Y. Bugeaud and M. Laurent,
Transcendence and continued fraction expansion of values of Hecke--Mahler series,
\emph{Acta Arith.} \textbf{209} (2023), 59--90.
\href{https://doi.org/10.4064/aa220323-18-1}{DOI}.

\bibitem{mMahler} B. Adamczewski and C. Faverjon, Mahler's method in several variables and
finite automata, \emph{Ann. of Math.} \textbf{204} (2026), 455--533.
\href{https://doi.org/10.4007/annals.2026.204.2.1}{DOI}.

\bibitem{mBhowmik} T. Bhowmik and S.~S. Pathak, A note on transcendence of special values
of functions related to modularity, \emph{Acta Arith.} \textbf{218} (2025), 1--23.
\href{https://doi.org/10.4064/aa240128-7-9}{DOI}.

\bibitem{mDiamond} F. Diamond and J. Shurman, \emph{A First Course in Modular Forms},
Graduate Texts in Mathematics 228, Springer, 2005.
\href{https://doi.org/10.1007/978-0-387-27226-9}{DOI}.

\bibitem{mMednykh} A.~D. Mednykh, On hyperbolic and spherical volumes for knot and link
cone-manifolds, in \emph{Kleinian Groups and Hyperbolic 3-Manifolds}, Cambridge
University Press, 2003, 145--164.
\href{https://doi.org/10.1017/CBO9780511542817.008}{DOI}.

\bibitem{mCDT} F. Calegari, V. Dimitrov and Y. Tang,
The linear independence of $1$, $\zeta(2)$, and $L(2,\chi_{-3})$,
arXiv:2408.15403, 2024. \href{https://arxiv.org/abs/2408.15403}{Preprint}.

\bibitem{mTang} Qian Tang, A Lean proof of the linear independence of
$1,\zeta(2),\zeta(3)$, public Lean 4 project, 2026. The theorem
\texttt{one\_zeta\_two\_zeta\_three\_linearIndependent} is in
\href{https://github.com/dtq1997/zeta2-zeta3-linear-independence}{the project repository}.

\bibitem{nCG} T.~Chatterjee and S.~Garg, \emph{On arithmetic nature of
$q$-analogue of the generalized Stieltjes constants}, 2024, Eq.~(7).
\href{https://arxiv.org/abs/2404.08025}{arXiv:2404.08025}.

\bibitem{nDKP} A.~B.~Dixit, V.~Kumar and S.~S.~Pathak,
Linear independence of values of the $q$-exponential and related functions,
\emph{Bull. Aust. Math. Soc.} \textbf{109} (2024), no.~3.
\href{https://doi.org/10.1017/S0004972723001028}{doi:10.1017/S0004972723001028}.

\bibitem{nTach} Y.~Tachiya, Irrationality of certain Lambert series,
\emph{Tokyo J. Math.} \textbf{27} (2004), 75--85.
\href{https://doi.org/10.3836/tjm/1244208475}{doi:10.3836/tjm/1244208475}.
See also the author's
\href{https://www.kurims.kyoto-u.ac.jp/~kyodo/kokyuroku/contents/pdf/1384-34.pdf}{RIMS exposition},
\emph{K\=oky\=uroku} \textbf{1384} (2004), 231--238.

\bibitem{nBeuk} F.~Beukers, A note on the irrationality of $\zeta(2)$
and $\zeta(3)$, \emph{Bull. London Math. Soc.} \textbf{11} (1979),
268--272. \href{https://doi.org/10.1112/blms/11.3.268}{doi:10.1112/blms/11.3.268}.

\bibitem{nShid} A.~B.~Shidlovskii, \emph{Transcendental Numbers},
de Gruyter Studies in Mathematics 12, Berlin, 1989, Ch.~7, \S1,
Thm.~1. For the normalization $w_j(z)$ and the application at $-1$, see
M.~R.~Powers, \emph{Arithmetic properties of several generalized-constant
sequences, with implications for $\Gamma^{(n)}(1)$}, 2025,
\href{https://arxiv.org/abs/2511.01849}{arXiv:2511.01849}.

\bibitem{nDST} D.~Duverney, Y.~Suzuki and Y.~Tachiya, \emph{Linear
independence results for certain sums of reciprocals of Fibonacci and
Lucas numbers}, 2019, Cor.~1.2 and Ex.~1.3.
\href{https://arxiv.org/abs/1908.07290}{arXiv:1908.07290}.

\bibitem{nKeio} Keio research archive, \emph{A Certificate-Driven
Continued-Fraction Constant}, November 11, 2026, private research
manuscript. The defining procedure is reproduced in the present paper.

\bibitem{nLamDraft} Keio and Shaitetsu, \emph{A Larger Irrationality
Region for Rational Lambert Bases}, research draft for independent
review, November 13, 2026.

\bibitem{nLamCode} Keio research archive,
\emph{Erdos1049\_Cutoff\_Check.py}, standalone rational-arithmetic
certificate, November 13, 2026.

\bibitem{nLamSource} C.~Lambrinoudis, \emph{Nonquadraticity of Lambert
Series Values}, 2026, manuscript and formalization.
\href{https://github.com/clambro/lambert-nonquadraticity/tree/008e3136a95e1e1ea605837ca9742b8613b2034a}{Repository, revision 008e3136a95e}.

\bibitem{pErdos} P. Erd\H{o}s, On the irrationality of certain series: problems and results (1988). \href{https://renyi.hu/~p_erdos/1988-22.pdf}{Author's paper}. See also T. Bloom, \href{https://www.erdosproblems.com/68}{Erd\H{o}s Problem 68}.
\end{thebibliography}
\end{document}
